\documentclass[11pt,a4paper]{article}
\usepackage[utf8]{inputenc}
\usepackage[T1]{fontenc}
\usepackage[czech]{babel}
\usepackage{lmodern}
\usepackage{geometry}
\usepackage{xcolor}
\usepackage{array}
\usepackage{longtable}
\usepackage{booktabs}
\usepackage{enumitem}
\usepackage{hyperref}
\usepackage{titlesec}

\geometry{margin=2.35cm}
\hypersetup{
  colorlinks=true,
  linkcolor=blue!55!black,
  urlcolor=blue!55!black,
  pdftitle={B-CS statnice - pravo a politika pro kyberbezpecnost},
  pdfauthor={Petr Marinec}
}
\setlist[itemize]{leftmargin=*, topsep=2pt, itemsep=2pt}
\setlist[enumerate]{leftmargin=*, topsep=2pt, itemsep=2pt}
\titleformat{\section}{\Large\bfseries}{\thesection}{0.8em}{}
\titleformat{\subsection}{\large\bfseries}{\thesubsection}{0.8em}{}
\titleformat{\subsubsection}{\normalsize\bfseries}{\thesubsubsection}{0.8em}{}
\newcommand{\otazka}[3]{\section*{#1}\addcontentsline{toc}{section}{#1}\noindent\textbf{Znění otázky:} #2\par\medskip}
\newcommand{\zkouska}[1]{\medskip\noindent\colorbox{blue!7}{\parbox{\dimexpr\linewidth-2\fboxsep}{\textbf{Co říct u zkoušky:} #1}}\medskip}
\newcommand{\pozor}[1]{\medskip\noindent\colorbox{red!7}{\parbox{\dimexpr\linewidth-2\fboxsep}{\textbf{Pozor:} #1}}\medskip}
\newcommand{\priklad}[1]{\medskip\noindent\colorbox{green!7}{\parbox{\dimexpr\linewidth-2\fboxsep}{\textbf{Příklad:} #1}}\medskip}
\newcommand{\pravobox}[1]{\medskip\noindent\colorbox{yellow!18}{\parbox{\dimexpr\linewidth-2\fboxsep}{\textbf{Právní kontext:} #1}}\medskip}
\newcommand{\zdroje}[1]{\medskip\noindent\textbf{Zdroje k opakování:} #1\medskip}

\title{B-CS státnice: právo a politika pro kyberbezpečnost\\\large Rozšířené učební poznámky k otázkám 15-23}
\author{248689}
\date{Aktualizováno 16. června 2026}

\begin{document}
\maketitle
\tableofcontents
\newpage

\section*{Jak s materiálem pracovat}
\addcontentsline{toc}{section}{Jak s materiálem pracovat}
Tyto poznámky vycházejí z přiloženého PDF \emph{Statnice B-CS blok 3, otázky 15-23, v4.2}, zachovávají jeho právní korekce a doplňují je do podoby širších studijních zápisků. U prvních tří BSS otázek jsou navíc zapracované dodané starší poznámky: širší vymezení bezpečnostních studií, vztah bezpečnostních a strategických studií, sektory, pojmy hrozba/riziko a základy kyberválky. U právních otázek je hlavním východiskem ověřený aktuální stav, ne starší právo.

U každé otázky je nejdříve kostra odpovědi, potom rozbor přesně podle podotázek ze zadání. U zkoušky je dobré nemluvit jen seznamem pojmů, ale vždy ukázat vztah: \emph{co chráním, před kým nebo čím, jaký je právní/institucionální režim, jaký je praktický dopad a kde je typická zkoušková chyba}.

\section*{Stavová tabulka k právu a strategiím}
\addcontentsline{toc}{section}{Stavová tabulka k právu a strategiím}
\begin{longtable}{p{0.25\linewidth}p{0.24\linewidth}p{0.20\linewidth}p{0.24\linewidth}}
\toprule
\textbf{Předpis / dokument} & \textbf{Stav k 16. 6. 2026} & \textbf{Použitelnost / účinnost} & \textbf{Co si pamatovat}\\
\midrule
NIS2 - směrnice 2022/2555 & Platná; transpoziční lhůta skončila 17. 10. 2024. & V ČR transponována hlavně zákonem o kybernetické bezpečnosti (č. 264/2025 Sb.). & Směrnice se neuplatní jako běžný český zákon sama o sobě; do praxe ji převádí národní úprava.\\
\addlinespace
Zákon o kybernetické bezpečnosti (č. 264/2025 Sb.) & Účinný, nahradil starý zákon o kybernetické bezpečnosti z roku 2014. & 1. 11. 2025. & Režim vyšších a nižších povinností, sebeidentifikace regulované služby, hlášení incidentů, odpovědnost vedení.\\
\addlinespace
Vyhlášky k novému ZKB & Účinné od 1. 11. 2025. & Zejména vyhláška o regulovaných službách, vyhláška o bezpečnostních opatřeních pro vyšší režim a vyhláška o bezpečnostních opatřeních pro nižší režim. & U zkoušky je důležité rozlišit, jestli subjekt vůbec poskytuje regulovanou službu a v jakém režimu.\\
\addlinespace
Národní strategie kybernetické bezpečnosti 2026-2030 & Zveřejněna NÚKIB 8. 9. 2025. & Strategický dokument od roku 2026. & Navazuje na Bezpečnostní strategii ČR 2023, Obrannou strategii 2023, EU a NATO dokumenty; akční plán je samostatný implementační dokument.\\
\addlinespace
CER - směrnice 2022/2557 a zákon o odolnosti subjektů kritické infrastruktury (č. 266/2025 Sb.) & Nová úprava odolnosti kritických subjektů. & Zákon účinný od 19. 8. 2025, praktické zařazování probíhá v přechodných lhůtách. & Vedle kybernetické bezpečnosti řeší fyzickou a organizační odolnost kritických subjektů.\\
\addlinespace
DORA - nařízení 2022/2554 & Platné a použitelné. & 17. 1. 2025. & Lex specialis pro finanční sektor: ICT rizika, incidenty, testování digitální odolnosti, dohled nad ICT třetími stranami.\\
\addlinespace
CRA - nařízení 2024/2847 & V platnosti. & Hlášení od 11. 9. 2026, hlavní povinnosti od 11. 12. 2027. & Produkty s digitálními prvky: bezpečný vývoj, zranitelnosti, aktualizace, CE logika.\\
\addlinespace
EU e-evidence - nařízení 2023/1543 & Platné. & Použitelné od 18. 8. 2026. & EPOC je evropský příkaz k vydání, EPOC-PR je příkaz k uchování. Neplést lhůty vydání a uchování.\\
\addlinespace
EU e-evidence - směrnice 2023/1544 & Platná; členské státy měly transponovat do 18. 2. 2026. & Navazuje na použitelnost e-evidence režimu. & Určené provozovny a právní zástupci poskytovatelů pro příjem příkazů.\\
\addlinespace
eIDAS 2.0 - nařízení 2024/1183 & V platnosti od 20. 5. 2024. & EUDI peněženky se odvíjejí od prováděcích aktů; obecně se uvádí horizont roku 2026, ale přesné lhůty je třeba ověřit podle aktuálních aktů. & Kvalifikovaný elektronický podpis má účinky vlastnoručního podpisu; EUDI peněženka přidává digitální identitu.\\
\addlinespace
2. dodatkový protokol k Budapešťské úmluvě - CETS 224 & Stále není v platnosti, protože je potřeba pět ratifikací. & Vstup po splnění ratifikační podmínky. & Nelze ho uvádět jako už účinný právní základ; je to budoucí nástroj přímější spolupráce.\\
\addlinespace
Úmluva OSN proti kyberkriminalitě & Není v platnosti; UNTC uvádí 77 signatářů a 3 strany. & Vstup 90 dní po 40. ratifikaci. & Podpis není ratifikace. U zkoušky říct, že jde zatím hlavně o politický a budoucí kooperační rámec.\\
\bottomrule
\end{longtable}

\newpage

\otazka{15. Vymezení bezpečnostních studií}{Pojmy bezpečnost, hrozba a riziko a jejich použití v oblasti kyberbezpečnosti. Vnitřní a vnější bezpečnost a sektory z pohledu kyberbezpečnosti. Vymezení bezpečnostní politiky a její analýza se zaměřením na politiku v oblasti kyberbezpečnosti.}{BSSb1101}

\subsection*{Kostra odpovědi}
\begin{enumerate}
  \item Vymezit bezpečnostní studia: co zkoumají, proč se po studené válce rozšířila a jak souvisejí se strategickými studii.
  \item Vysvětlit řetězec bezpečnost - hrozba - zranitelnost - riziko, vždy obecně i kyberneticky.
  \item Ukázat rozdíl mezi vnitřní a vnější bezpečností a vysvětlit, proč ho kyberprostor rozmazává.
  \item Popsat sektory bezpečnosti podle Kodaňské školy a pojem sekuritizace.
  \item Vymezit bezpečnostní politiku a nabídnout postup její analýzy v kyberbezpečnosti.
\end{enumerate}

\subsection*{Bezpečnostní studia jako obor}
Bezpečnostní studia zkoumají, jaké hodnoty jsou ohroženy, odkud ohrožení přichází, kdo je za bezpečnost odpovědný a jakými nástroji lze bezpečnost zajišťovat. V klasickém pojetí byla bezpečnost spojena hlavně se státem, armádou a válkou. Proto se bezpečnostní studia dlouho brala jako subdisciplína mezinárodních vztahů a strategická studia byla jejich vojensky orientovanou částí.

Po studené válce se pojetí bezpečnosti rozšířilo. Do centra se dostala nejen vojenská agrese, ale i ekonomická závislost, terorismus, organizovaný zločin, energetika, migrace, pandemie, klimatické dopady, informační operace a kybernetické útoky. Po 11. září 2001 se významně posílila také studia vnitřní bezpečnosti, protože hrozby už nebylo možné pohodlně rozdělit na "venku je armáda" a "uvnitř je policie".

Bezpečnostní studia jsou proto průřezová. Berou si pojmy z politologie, mezinárodních vztahů, práva, sociologie, ekonomie, psychologie, historie, kriminologie, vojenské vědy a technických oborů. Kyberbezpečnost je v tom typická: nejde čistě o techniku, protože útok na nemocnici je zároveň technický incident, právní problém, krizové řízení, otázka lidské bezpečnosti, ekonomická škoda a často i zahraničněpolitický signál.

Když to máš říct nahlas, neříkej jen "bezpečnostní studia řeší bezpečnost". Lepší je postup: nejdřív pojmenuj chráněnou hodnotu, potom hrozbu, potom aktéry a nakonec reakci. Například u útoku na nemocnici chráněnou hodnotou není jen nemocniční server, ale poskytování zdravotní péče, život a zdraví pacientů, důvěrnost zdravotnických dat a důvěra veřejnosti. Hrozbou může být ransomwarová skupina nebo státem podporovaný aktér. Aktéři obrany nejsou jen IT správci, ale vedení nemocnice, dodavatel systému, NÚKIB, policie, ÚOOÚ, zřizovatel a někdy i mezinárodní partneři. Reakce není jen "nahodit zálohu", ale technická obnova, trestní řízení, hlášení incidentu, krizová komunikace a poučení do budoucna.

\zkouska{Začni větou: "Bezpečnostní studia nejsou jen studium války, ale širší obor zkoumající hrozby vůči chráněným hodnotám a způsoby jejich zvládání." Potom hned ukaž, že kyberbezpečnost je dobrý příklad rozšířeného pojetí bezpečnosti.}

\subsection*{Bezpečnost, hrozba, zranitelnost a riziko}
\subsubsection*{Bezpečnost}
Bezpečnost je stav i proces. Jako stav znamená, že chráněná hodnota není vystavena nepřijatelnému nebezpečí. Jako proces znamená soustavné snižování hrozeb, zranitelností a dopadů na přijatelnou úroveň. Důležité je ptát se: bezpečnost koho nebo čeho? Referenčním objektem může být stát, obyvatelstvo, jednotlivec, firma, nemocnice, informační systém, data nebo důvěra ve volby.

Rozlišuje se také \textbf{objektivní} a \textbf{subjektivní} bezpečnost. Objektivní bezpečnost se ptá, jaké hrozby a zranitelnosti skutečně existují a jaká je pravděpodobnost a dopad škody. Subjektivní bezpečnost je pocit bezpečí nebo ohrožení. V kyberprostoru se mohou rozcházet: firma může mít objektivně slabé zabezpečení, ale vedení má falešný pocit bezpečí, protože "se nám ještě nic nestalo". Naopak po mediálně viditelném incidentu může veřejnost vnímat digitální služby jako nebezpečné, i když konkrétní riziko bylo rychle omezeno.

V kyberbezpečnosti se bezpečnost často konkretizuje přes triádu CIA:
\begin{itemize}
  \item \textbf{Důvěrnost} - k datům se dostane jen oprávněná osoba.
  \item \textbf{Integrita} - data a systémy nejsou neoprávněně změněny.
  \item \textbf{Dostupnost} - služba funguje, když ji uživatel potřebuje.
\end{itemize}
U kritických služeb je potřeba doplnit i odolnost, schopnost obnovy, auditovatelnost, autenticitu a bezpečnost dodavatelského řetězce.

Je také užitečné rozlišit \emph{security} a \emph{safety}. Security míří na úmyslné jednání, například hacking, sabotáž nebo špionáž. Safety míří na neúmyslné nebo přírodní události, například požár, povodeň, technické selhání nebo chybu obsluhy. V praxi se prolínají: výpadek nemocničního systému může začít jako security incident, ale dopady má na safety pacientů.

\textbf{Bezpečnostní dilema} znamená situaci, kdy opatření jednoho aktéra ke zvýšení vlastní bezpečnosti snižuje pocit bezpečí druhého aktéra, a ten proto posiluje vlastní kapacity. Výsledkem může být spirála nedůvěry. V kyberprostoru je to časté: stát buduje ofenzivní kyberkapacity nebo hromadí informace o zranitelnostech, aby se uměl bránit a odstrašovat, ale jiné státy to mohou číst jako přípravu útoku. Podobně povinné monitorování sítí zvyšuje bezpečnost, ale může vyvolat obavy ze zásahu do soukromí. Proto se bezpečnostní politika musí ptát nejen "co nás ochrání", ale i "jaké vedlejší dopady a signály opatření vytváří".

\subsubsection*{Hrozba}
Hrozba je potenciální zdroj škody. Může být intencionální, tedy spojená s vůlí aktéra, nebo neintencionální, například přírodní událost či lidská chyba. V kyberbezpečnosti je hrozbou ransomwarová skupina, státem podporovaná APT skupina, insider, botnet, chyba administrátora, požár serverovny nebo výpadek dodavatele cloudu.

Hrozby lze dál popsat jako aktivní, pasivní, pokročilé nebo trvalé:
\begin{itemize}
  \item aktivní hrozba se snaží přímo škodit, typicky malware, DDoS nebo kompromitace účtu,
  \item pasivní hrozba sleduje, odposlouchává nebo sbírá informace,
  \item pokročilá hrozba používá sofistikované metody, zero-day zranitelnosti a cílený spear phishing,
  \item trvalá hrozba odpovídá APT logice - dlouhodobé a skryté působení v systému.
\end{itemize}

\subsubsection*{Zranitelnost}
Zranitelnost je slabina, přes kterou hrozba může způsobit škodu. Technicky to může být neopravená chyba, špatná konfigurace, slabé heslo nebo chybějící MFA. Organizačně to může být chybějící plán obnovy, nejasné odpovědnosti, neproškolení lidé nebo závislost na jednom dodavateli. Zranitelnost je prakticky nejlépe ovlivnitelná: opravy, segmentace, zálohy, řízení přístupů, školení, monitoring a testování obnovy.

\subsubsection*{Riziko}
Riziko je vztah hrozby ke konkrétnímu referenčnímu objektu. Jednoduše se říká, že riziko je pravděpodobnost krát dopad, ale je to jen pomůcka. Přesnější je vnímat riziko jako kombinaci pravděpodobnosti, dopadu, zranitelnosti, expozice, nejistoty a schopnosti reakce. Stejná hrozba vytváří jiné riziko pro malý blog, jiné pro banku a jiné pro nemocnici.

\priklad{Ransomwarová skupina je hrozba. Nemocnice se starými systémy, bez segmentace a bez ověřených offline záloh má vysoké riziko, protože dopad na péči bude kritický. Firma se segmentací, EDR, testovanými zálohami a plánem obnovy čelí stejné hrozbě, ale riziko je nižší.}

\pozor{Neříkej "hrozba rovná se riziko". Hrozba je zdroj možného poškození, riziko je míra ohrožení konkrétního objektu.}

\subsection*{Vnitřní a vnější bezpečnost v kyberprostoru}
Vnější bezpečnost tradičně chrání stát před hrozbami zvenčí: vojenskou agresí, nepřátelskými státy, zahraničními zpravodajskými službami nebo mezinárodním nátlakem. Nástroje jsou diplomacie, obrana, aliance a zpravodajství. Vnitřní bezpečnost řeší stabilitu a pořádek uvnitř státu: kriminalitu, terorismus, extremismus, ochranu obyvatelstva, veřejný pořádek a fungování institucí.

Kyberprostor toto dělení stírá ze tří důvodů. Zaprvé je neteritoriální: útok může být řízen z jiného kontinentu, infrastruktura může být v cloudu ve třetí zemi a dopad se projeví v české nemocnici. Zadruhé je atribuce nejistá: nevíme hned, jestli útočí stát, kriminální skupina, hacktivisté nebo proxy aktér. Zatřetí většinu relevantní infrastruktury vlastní soukromý sektor, takže bezpečnost nelze zajistit jen státním rozkazem.

Proto se v kyberbezpečnosti potkává NÚKIB, policie, zpravodajské služby, armáda, ministerstva, kraje, obce, regulované organizace, poskytovatelé služeb, dodavatelé technologií, CERT/CSIRT týmy a mezinárodní partneři. Kybernetický incident může být současně trestným činem, porušením GDPR, incidentem podle ZKB, krizovou situací a zahraničněpolitickou otázkou.

Plynulé vysvětlení u zkoušky může vypadat takto: "V klasickém pojetí by vnější bezpečnost řešila cizí stát a armádu, zatímco vnitřní bezpečnost policii a kriminalitu. Kyberprostor to komplikuje, protože stejný incident může začít jako phishing zaměstnance v české nemocnici, technicky běžet přes servery ve třetí zemi, být financovaný kriminální skupinou a současně být tolerovaný nepřátelským státem. Proto se nedá říct, že kyberbezpečnost je jen vnitřní nebo jen vnější bezpečnost. Je to průřezová agenda, kde se potkává policie, zpravodajství, diplomacie, regulace i soukromý sektor."

\subsection*{Sektory bezpečnosti a sekuritizace}
Kodaňská škola rozšířila bezpečnost do pěti sektorů:
\begin{itemize}
  \item \textbf{vojenský} - obrana, ozbrojené síly, schopnost odstrašení,
  \item \textbf{politický} - suverenita, legitimita, stabilita institucí,
  \item \textbf{ekonomický} - finance, obchod, energie, dodavatelské řetězce,
  \item \textbf{societální} - identita, sociální soudržnost, důvěra,
  \item \textbf{environmentální} - životní prostředí a přírodní zdroje.
\end{itemize}
Kyberbezpečnost je průřezová. Útok na energetiku je vojenský i ekonomický problém, ransomware v nemocnici je lidská bezpečnost a kritická infrastruktura, dezinformační operace míří na politický a societální sektor.

Sekuritizace je proces, kdy aktér prohlásí téma za existenční hrozbu pro referenční objekt a publikum toto rámování přijme. Nestačí tedy, že politik řekne "tohle je bezpečnostní hrozba". Úspěšná sekuritizace nastává až tehdy, když relevantní publikum přijme mimořádná opatření. Kyberbezpečnost se za poslední dvě dekády sekuritizovala: z technické otázky správců sítí se stala součást národní bezpečnosti, vznikl samostatný úřad, zákon, strategie, povinnosti organizací a mezinárodní sankční/diplomatická reakce.

Celé schéma sekuritizace je dobré umět přesně. \textbf{Sekuritizující aktér} je ten, kdo hrozbu pojmenovává, například vláda, NÚKIB, ministr, zpravodajská služba nebo vedení organizace. \textbf{Referenční objekt} je to, co má být existenčně ohroženo: stát, obyvatelstvo, nemocniční péče, volby, finanční stabilita nebo důvěra ve veřejnou správu. \textbf{Existenční hrozba} je rámování, že nejde o běžný provozní problém, ale o ohrožení základního fungování. \textbf{Publikum} jsou ti, kdo musí rámování přijmout: Parlament, vláda, regulované subjekty, veřejnost, média nebo profesní komunita. \textbf{Funkční aktéři} nejsou hlavní mluvčí, ale ovlivňují výsledek, například dodavatelé cloudu, operátoři, nemocnice, banky, experti nebo média. \textbf{Mimořádná opatření} jsou kroky, které by bez bezpečnostního rámování nebyly politicky nebo právně přijatelné: nové povinnosti, mimořádné rozpočty, varování, zákaz rizikového dodavatele, krizová opatření nebo rozsáhlejší monitoring.

Kodaňská škola zároveň rozlišuje stupně: \textbf{nepolitizované} téma není ve veřejné politice skoro řešeno, \textbf{politizované} téma je předmětem běžného rozhodování a rozpočtů, \textbf{sekuritizované} téma je prezentováno jako existenční hrozba vyžadující mimořádná opatření. \textbf{Desekuritizace} je opačný proces: téma se vrací z výjimečného bezpečnostního režimu do běžné politiky. U kyberbezpečnosti je cílem často vyváženost: některé incidenty je nutné sekuritizovat, ale běžná správa IT nemá být trvale vedena v režimu permanentní paniky.

Sektory není potřeba jen vyjmenovat; je potřeba ukázat, že jeden kyberincident může přeskakovat mezi nimi. Útok na distribuční síť elektřiny je ekonomický, protože ohrožuje výrobu a služby, politický, protože testuje schopnost státu zajistit základní funkce, vojenský, pokud je součástí nátlaku nepřátelského státu, a societální, pokud vyvolá strach a nedůvěru. Právě proto se kyberbezpečnost často označuje jako průřezová bezpečnostní agenda.

\priklad{Sekuritizace v praxi: po útocích na české nemocnice se z problému "nemocnice mají slabé IT" stal problém "výpadek nemocničního IT ohrožuje životy a fungování státu". Jakmile toto rámování přijali politici, veřejnost a instituce, následovala mimořádnější opatření: metodiky, finance, povinnosti, cvičení, větší role NÚKIB a tlak na odolnost zdravotnictví.}

\subsection*{Bezpečnostní politika a její analýza}
Bezpečnostní politika je nejobecnější program státu pro ochranu základních hodnot: suverenity, územní celistvosti, demokratického řádu, života a zdraví obyvatel, fungování institucí, ekonomiky a mezinárodních závazků. Není to jen jeden dokument. Je to celý způsob, jak stát odpovídá na otázky: co chráníme, před čím to chráníme, kdo za to odpovídá, jaké prostředky použijeme a jak poznáme, že to funguje.

V praxi se bezpečnostní politika skládá z několika propojených částí. Zahraniční bezpečnostní politika řeší spojence, diplomacii, sankce a mezinárodní normy. Obranná politika řeší armádu, odstrašení a schopnost bránit stát. Politika vnitřní bezpečnosti řeší policii, kriminalitu, terorismus, krizové řízení a ochranu obyvatelstva. Hospodářská bezpečnostní politika řeší energetiku, dodavatelské řetězce, strategické suroviny a odolnost ekonomiky. Informační a kybernetická politika řeší data, sítě, digitální služby, dezinformace, důvěru a bezpečnost technologií.

Na kyberbezpečnosti je dobře vidět, že bezpečnostní politika není jen "stát vydá zákaz". Stát často nevlastní nemocnice, bankovní infrastrukturu, cloudové služby ani telekomunikační sítě. Musí proto kombinovat regulaci, metodickou podporu, dohled, financování, vzdělávání, sdílení informací, policejní vyšetřování, zpravodajské varování, diplomacii a mezinárodní spolupráci. Kyberpolitika je tedy zároveň právní, technická, organizační, ekonomická i zahraničněpolitická.

U analýzy je dobré rozlišit tři roviny politiky. \textbf{Polity} je institucionální rámec: kdo má kompetence, jaké jsou úřady, zákony a vztahy mezi nimi. V kyberbezpečnosti sem patří NÚKIB, vláda, BRS, policie, zpravodajské služby, armáda, ÚOOÚ, ČTÚ, DIA a regulované subjekty. \textbf{Politics} je proces rozhodování: kdo prosazuje přísnější regulaci, kdo se bojí nákladů, jak se vyjednává mezi státem a soukromým sektorem, jak do toho vstupuje EU, NATO nebo veřejné mínění. \textbf{Policy} je konkrétní obsah: zákon o kybernetické bezpečnosti, bezpečnostní opatření, hlášení incidentů, strategie, akční plány, cvičení a dotace.

Analýzu bezpečnostní politiky je dobré dělat podle jednoduché šablony:
\begin{enumerate}
  \item \textbf{Referenční objekt} - co chráníme: stát, službu, data, volby, pacienty?
  \item \textbf{Hrozby a rizika} - kdo nebo co škodí, jaké jsou zranitelnosti a dopady?
  \item \textbf{Aktéři} - kdo rozhoduje, kdo provozuje, kdo útočí, kdo pomáhá?
  \item \textbf{Nástroje} - právo, strategie, regulace, technická opatření, diplomacie, trestní právo, vzdělávání.
  \item \textbf{Priority a trade-offy} - kolik to stojí, jaký je dopad na práva, podnikání a inovace?
  \item \textbf{Vyhodnocení} - jak poznáme, že politika funguje?
\end{enumerate}

Na příkladu nemocnic: referenční objekt nejsou jen servery, ale dostupnost zdravotní péče, pacientská data a důvěra lidí, že nemocnice bude fungovat. Hrozby jsou ransomware, phishing, zneužití dodavatele, insider nebo výpadek cloudu. Zranitelnosti mohou být staré systémy, slabá segmentace, chybějící offline zálohy a neproškolený personál. Aktéři jsou nemocnice, zřizovatel, dodavatelé, NÚKIB, policie, ÚOOÚ, ministerstvo zdravotnictví, kraj a případně zahraniční partneři. Nástroje jsou regulace, bezpečnostní opatření, cvičení, incident reporting, trestní řízení, metodiky, financování a krizová komunikace. Trade-off je třeba cena bezpečnosti, dopad na provoz nemocnice a ochrana soukromí pacientů. Vyhodnocení se nedělá podle pocitu, ale podle toho, zda nemocnice umí incident včas zjistit, omezit, nahlásit, obnovit provoz a snížit dopad na pacienty.

\priklad{Když stát zavede povinnost hlásit významné kyberincidenty, nejde jen o byrokracii. Politický cíl je, aby stát rychle viděl systémovou hrozbu, mohl varovat další subjekty a koordinovat reakci. Současně ale musí řešit proporcionalitu: hlášení nesmí být tak složité, aby zahltilo nemocnici uprostřed krize, a stát musí chránit citlivé informace, které od subjektu získá.}

\zkouska{Když dostaneš otázku na analýzu kyberpolitiky, mluv jako analytik: "Nejdřív řeknu, co přesně chráníme. Pak určím hrozby, zranitelnosti a dopady. Potom pojmenuji aktéry a jejich odpovědnosti. Následně vysvětlím nástroje - právo, instituce, technická opatření, peníze, diplomacii a trestní právo. Nakonec zhodnotím, jestli jsou opatření přiměřená, proveditelná a měřitelná."}

\zdroje{Původní PDF v4.2; dodané BSS poznámky; NIS2 definice v čl. 6 směrnice 2022/2555; Buzan, Waever, de Wilde: \emph{Security: A New Framework for Analysis}.}

\newpage

\otazka{16. Bezpečnostní strategie}{Dokumenty České republiky a význam kyberbezpečnosti v nich. Bezpečnostní systém České republiky a role institucí v oblasti kyberbezpečnosti.}{BSSb1103}

\subsection*{Kostra odpovědi}
\begin{enumerate}
  \item Vysvětlit, co je bezpečnostní strategie a jak se liší od zákona a akčního plánu.
  \item Vyjmenovat klíčové dokumenty ČR a zařadit kyberbezpečnost.
  \item Popsat bezpečnostní systém ČR: prezident, Parlament, vláda, BRS, ÚKŠ, ministerstva, kraje, složky.
  \item Vysvětlit roli NÚKIB, MV, MO/Armády, policie, zpravodajských služeb, DIA, ČTÚ a soukromého sektoru.
  \item Ukázat, jak strategie vede k právu, institucionálnímu systému a praktické reakci na incident.
\end{enumerate}

\subsection*{Co je bezpečnostní strategie}
Bezpečnostní strategie je vrcholný politicko-strategický dokument. Nezakládá sama o sobě jednotlivci povinnosti jako zákon, ale určuje, jak stát chápe bezpečnostní prostředí, jaké má zájmy, jaké hrozby považuje za zásadní a jakými nástroji chce reagovat. Je to most mezi analýzou hrozeb a konkrétní politikou: z ní se odvozují rezortní strategie, akční plány, rozpočty, organizační změny a legislativa. U zkoušky je dobré ji vysvětlit jako "mapu a zadání": strategie řekne, kam stát míří a proč; zákony, vyhlášky, instituce a rozpočty potom řeší, kdo to udělá a jak.

Rozdíl mezi strategiemi je v úrovni obecnosti. Bezpečnostní strategie ČR je nejširší rámec. Obranná strategie řeší obranu státu a vojenskou dimenzi. Národní strategie kybernetické bezpečnosti řeší digitální a kybernetickou část bezpečnosti. Akční plán je implementační dokument: překládá strategické cíle do úkolů, termínů a odpovědných institucí. Koncepční dokument tedy obvykle neříká "tato nemocnice musí mít přesně toto MFA", ale říká "zvýšíme odolnost zdravotnictví a kritických služeb"; konkrétní povinnosti se pak objeví v zákoně, vyhlášce, metodice nebo dotačním programu.

\subsection*{Dokumenty ČR a kyberbezpečnost}
\textbf{Bezpečnostní strategie ČR 2023} vychází z výrazně zhoršeného bezpečnostního prostředí po ruské agresi proti Ukrajině. Zdůrazňuje, že hrozby jsou propojené: vojenské, hybridní, kybernetické, ekonomické, energetické a informační. Kyberbezpečnost v ní není okrajová IT agenda, ale součást ochrany suverenity, demokracie, kritické infrastruktury, ekonomiky a důvěry v instituce.

\textbf{Obranná strategie ČR 2023} zdůrazňuje obranyschopnost, odstrašení, vazbu na NATO a potřebu připravit se i na operace v kyberprostoru a informačním prostředí. Pro státnice je důležité říct, že kyberobrana není totéž co kybernetická bezpečnost: kyberobrana je vojenská dimenze a spadá hlavně pod Ministerstvo obrany a Armádu ČR, zatímco civilní kyberbezpečnost koordinuje NÚKIB.

\textbf{Národní strategie kybernetické bezpečnosti 2026-2030} je aktuální specializovaný dokument. NÚKIB ji zveřejnil v roce 2025 jako strategii platnou pro období 2026-2030. Navazuje na Bezpečnostní strategii ČR 2023, Obrannou strategii 2023, Strategickou koncepci NATO 2022 a unijní kybernetickou strategii. Důraz klade na odolnost státu a kritických služeb, bezpečnost dodavatelských řetězců, technologickou suverenitu, sdílení informací, detekci a reakci, lidské zdroje, mezinárodní spolupráci a koordinované odstrašení škodlivých aktérů.

Strategie 2026-2030 se dá u zkoušky dobře zapamatovat přes tři velké směry. První je \textbf{bezpečná strategická infrastruktura}: chránit základní služby, zvyšovat odolnost státu, umět včas detekovat a zvládat incidenty, řešit dodavatelské řetězce a stavět veřejné IT tak, aby bezpečnost nebyla až dodatečná náplast. Druhý je \textbf{celospolečenská připravenost a rozvoj}: vzdělávání, nedostatek expertů, bezpečnostní kultura, výzkum, inovace a koordinace mezi státem, školami, firmami a odbornou komunitou. Třetí je \textbf{mezinárodní spolupráce a prosazování zájmů}: aktivní role v EU, NATO a dalších fórech, sdílení informací, společné atribuce, sankce, diplomatické reakce a podpora otevřeného, bezpečného a svobodného kyberprostoru.

\textbf{Další relevantní dokumenty} jsou například Koncepce rozvoje NÚKIB, Národní politika koordinovaného zveřejňování zranitelností, strategické dokumenty k hybridnímu působení, krizovému řízení, ochraně kritické infrastruktury a digitální transformaci veřejné správy. Pro kritickou infrastrukturu je nově důležitá také \textbf{Strategie pro posílení odolnosti subjektů kritické infrastruktury pro období 2026-2029}, kterou vláda schválila v prosinci 2025. Navazuje na CER logiku a dává politický rámec pro to, aby se kritické subjekty neposuzovaly jen podle kybernetických incidentů, ale podle all-hazards odolnosti základních služeb. U zkoušky není nutné vyjmenovat všechno, ale je dobré ukázat hierarchii: obecná bezpečnostní strategie - rezortní/specializovaná strategie - akční plán - zákony a vyhlášky - metodiky.

Plynulá odpověď může znít: "Strategické dokumenty nejsou právně závazné jako zákon, ale jsou důležité, protože nastavují politickou vůli a prioritu. Když Bezpečnostní strategie ČR řekne, že kyberprostor je součást národní bezpečnosti, není to jen prohlášení do šuplíku. Znamená to, že stát bude chtít posilovat NÚKIB, měnit zákony, vyčlenit peníze, cvičit instituce, řešit dodavatele a zapojovat se do EU a NATO. Právní norma pak z této politické priority udělá konkrétní povinnost."

\pozor{Nepoužívej starou formulaci, že strategie 2026-2030 není zveřejněná. Zveřejněná je; samostatný akční plán je potřeba ověřovat zvlášť.}

\zkouska{Jednoduchá věta k dokumentům: "Strategie pojmenovává bezpečnostní problém a cíle, akční plán rozděluje úkoly a termíny, zákon ukládá povinnosti a metodiky pomáhají subjektům povinnosti prakticky splnit."}

\subsection*{Bezpečnostní systém ČR}
Bezpečnostní systém ČR je soubor institucí, pravidel a procesů, jimiž stát zajišťuje bezpečnost. Ústavní rámec tvoří zejména Ústava, Listina, ústavní zákon o bezpečnosti České republiky (č. 110/1998 Sb.), krizová legislativa a zvláštní zákony.

Základní logika systému:
\begin{itemize}
  \item \textbf{Prezident republiky} má ústavní roli hlavy státu, jmenuje vládu, je vrchním velitelem ozbrojených sil a může se účastnit jednání Bezpečnostní rady státu; v kyberbezpečnosti ale není běžným řídícím orgánem civilní regulace.
  \item \textbf{Parlament} přijímá zákony, rozhoduje o některých bezpečnostních otázkách a kontroluje vládu.
  \item \textbf{Vláda} je hlavní výkonný orgán bezpečnostní politiky.
  \item \textbf{Bezpečnostní rada státu} je stálý pracovní orgán vlády pro koordinaci bezpečnostních otázek, nikoli samostatný ústavní orgán stojící vedle vlády; v jejím systému působí odborné výbory, v kybernetické oblasti zejména Výbor pro kybernetickou bezpečnost.
  \item \textbf{Ústřední krizový štáb} podporuje vládu při krizových situacích.
  \item \textbf{Ministerstva a jiné ústřední orgány} mají gesce podle oblastí: vnitro, obrana, zahraničí, průmysl, zdravotnictví, doprava atd.
  \item \textbf{Kraje, obce a složky IZS} zajišťují konkrétní výkon krizového řízení a ochrany obyvatelstva.
\end{itemize}

Kybernetická bezpečnost do toho vstupuje průřezově. Incident může vyžadovat technickou reakci, právní posouzení, krizové řízení, komunikaci s veřejností, trestní řízení a mezinárodní koordinaci. Proto nestačí jeden úřad: NÚKIB je gestor, ale není jediný aktér.

Bezpečnostní systém je vhodné vysvětlit jako řetězec od politického rozhodnutí k praktické reakci. Vláda určuje směr bezpečnostní politiky, Bezpečnostní rada státu připravuje koordinaci a návrhy opatření, odborné výbory řeší specializované agendy, ministerstva a úřady vykonávají své gesce a kraje, obce či složky IZS řeší konkrétní dopady v území. U kyberincidentu to znamená, že technický tým sice řeší škodlivý kód, ale zároveň někdo musí rozhodnout, zda je ohrožena základní služba, zda se aktivuje krizové řízení, zda se informuje veřejnost, zda se zapojí policie a zda se incident bude řešit i diplomaticky.

\subsection*{Krizové stavy a stav kybernetického nebezpečí}
Krizové stavy nejsou totéž co každý závažný incident. Závažný kybernetický útok může být důvodem ke krizovému řízení, ale sám automaticky neznamená vyhlášení nouzového stavu nebo válečného stavu. Rozhoduje rozsah dopadů, zda běžné nástroje nestačí a zda jsou splněny zákonné podmínky.

\begin{itemize}
  \item \textbf{Stav nebezpečí} vyhlašuje hejtman kraje nebo primátor hlavního města Prahy pro území kraje nebo jeho část, typicky při živelní pohromě, havárii nebo jiném ohrožení, které nelze zvládnout běžnými prostředky.
  \item \textbf{Nouzový stav} vyhlašuje vláda; při nebezpečí z prodlení ho může vyhlásit předseda vlády a vláda ho musí následně potvrdit nebo zrušit.
  \item \textbf{Stav ohrožení státu} vyhlašuje Parlament na návrh vlády, pokud je bezprostředně ohrožena svrchovanost, územní celistvost nebo demokratické základy státu.
  \item \textbf{Válečný stav} vyhlašuje Parlament, pokud je Česká republika napadena nebo je třeba plnit mezinárodní závazky o společné obraně.
\end{itemize}

Vedle toho existuje speciální \textbf{stav kybernetického nebezpečí} podle zákona o kybernetické bezpečnosti. Je to mimořádný nástroj pro závažné ohrožení kybernetické bezpečnosti, které nelze účinně zvládnout běžnými prostředky. Vyhlašuje ho ředitel NÚKIB nejdéle na 30 dnů; prodloužení je možné jen v zákonných mezích tak, aby celková doba nepřekročila 60 dnů. Nejde o ústavní krizový stav jako nouzový stav, ale o sektorový kybernetický režim umožňující rychlejší a silnější opatření vůči dotčeným systémům a subjektům.

\subsection*{Role institucí v kyberbezpečnosti}
\textbf{NÚKIB} je ústřední správní úřad pro kybernetickou bezpečnost. Vede regulaci podle zákona o kybernetické bezpečnosti, provozuje vládní CERT GovCERT.CZ, vydává varování, metodiky, opatření a koordinuje reakce na významné incidenty. Je také klíčovým aktérem strategické a mezinárodní kyberpolitiky.

\textbf{CSIRT.CZ} je národní CSIRT tým, který pomáhá hlavně s koordinací incidentů mimo vládní sektor. V praxi je důležité rozlišit vládní CERT a národní CSIRT.

\textbf{Policie ČR a státní zastupitelství} řeší kyberkriminalitu. Technickou expertizu mají specializované útvary, ale procesní rámec je trestní řád.

\textbf{Zpravodajské služby} sledují hrozby podle své působnosti: BIS vnitřní bezpečnost, ÚZSI zahraniční zpravodajství, Vojenské zpravodajství obrannou oblast. V kyberprostoru je jejich význam velký hlavně u státních aktérů, špionáže, atribuce a varování.

\textbf{Ministerstvo obrany a Armáda ČR} odpovídají za kyberobranu a vojenské působení v kyberprostoru. Národní centrum kybernetických operací a Velitelství kybernetických sil a informačních operací patří do obranné logiky státu. Je dobré je nezaměňovat s NÚKIB: NÚKIB je civilní regulátor a koordinátor kyberbezpečnosti, MO/Armáda řeší obranu a vojenské operace.

\textbf{Ministerstvo vnitra} je důležité pro krizové řízení, vnitřní bezpečnost, ochranu obyvatelstva, veřejnou správu a podle nové úpravy také odolnost kritických subjektů v nekybernetické dimenzi. \textbf{Ministerstvo zahraničních věcí} je důležité pro kyberdiplomacii, veřejnou atribuci, sankce a mezinárodní koordinaci. \textbf{DIA} je významná pro digitální služby státu, identitu občana a architekturu eGovernmentu. \textbf{ČTÚ} má roli v elektronických komunikacích a bude relevantní i u e-evidence režimu pro poskytovatele. \textbf{NBÚ} zůstává klíčový pro ochranu utajovaných informací, ale civilní kyberbezpečnost už není jeho hlavní gescí.

\textbf{Soukromý sektor} je nezastupitelný, protože vlastní nebo provozuje velkou část sítí, služeb a technologií. Stát proto kombinuje regulaci, dohled, sdílení informací, pobídky, metodickou podporu a krizovou koordinaci.

\priklad{Ransomware v nemocnici: nemocnice řeší obnovu, NÚKIB posuzuje incident podle ZKB a koordinuje, policie řeší trestný čin, ÚOOÚ řeší únik osobních údajů, kraj a zdravotnictví řeší kontinuitu péče, dodavatelé obnovují systémy a komunikační tým řeší důvěru veřejnosti.}

\priklad{Jak propojit dokumenty a instituce: Bezpečnostní strategie ČR pojmenuje kyberútoky jako hrozbu pro stát. Národní strategie kybernetické bezpečnosti z toho udělá cíle, například zvýšit odolnost strategické infrastruktury a vzdělávání. Zákon o kybernetické bezpečnosti určí regulované služby a povinnosti. NÚKIB vydává metodiky a dohlíží. Nemocnice nebo energetická firma pak musí zavést konkrétní opatření, hlásit incidenty a umět obnovit provoz.}

\zkouska{Dobrá odpověď propojí strategické dokumenty s institucemi: strategie říká, co je problém a cíl; zákony dávají pravomoci a povinnosti; instituce provádějí prevenci, detekci, reakci a obnovu.}

\zdroje{Bezpečnostní strategie ČR 2023; Obranná strategie ČR 2023; NÚKIB: Národní strategie kybernetické bezpečnosti 2026-2030; ústavní zákon o bezpečnosti České republiky (č. 110/1998 Sb.); zákon o kybernetické bezpečnosti (č. 264/2025 Sb.).}

\newpage

\otazka{17. Kyberválka}{Definice, historie, současné trendy. Identifikace aktérů, problémy s přisouzením kyberútoků, možnost odstrašení v kyberprostoru. Koncept netwars a jeho použití v kyberkonfliktu.}{BSSb1152}

\subsection*{Kostra odpovědi}
\begin{enumerate}
  \item Definovat kyberválku opatrně: ne každý kyberútok je válka.
  \item Popsat historii: od raných incidentů přes Estonsko 2007, Gruzii 2008, Stuxnet, Ukrajinu, NotPetya až po současné hybridní konflikty.
  \item Identifikovat aktéry: státy, proxy, kriminálníci, hacktivisté, firmy, jednotlivci.
  \item Vysvětlit problém atribuce: technická, operační, strategická a politická rovina.
  \item Vyložit odstrašení: by punishment, by denial, resilience, normy, sankce, právo.
  \item Definovat netwars a ukázat jejich vztah k informačním a kybernetickým konfliktům.
\end{enumerate}

\subsection*{Definice kyberválky}
Kyberválka je použití kybernetických prostředků v konfliktu mezi politickými aktéry, zejména státy, takovým způsobem, že účinky dosahují úrovně ozbrojeného konfliktu nebo jsou jeho součástí. Důležitá opatrnost: většina kyberútoků není kyberválka. Phishing, ransomware, špionáž nebo DDoS mohou být závažné, ale samy o sobě obvykle nepředstavují válku.

U státnic je nejlepší definovat kyberválku jako užší pojem než kyberkonflikt. Kyberkonflikt zahrnuje špionáž, sabotáž, informační operace, kriminální outsourcing a nátlak pod prahem ozbrojeného útoku. Kyberválka je až krajní forma, kdy kybernetické operace plní vojenské cíle, působí rozsáhlé fyzické nebo funkční škody, případně jsou integrovány do konvenční války.

Když to vysvětluješ, pomáhá začít rozdílem mezi válkou, kyberobranou, kyberkriminalitou a kyberbezpečností. Válka je politický konflikt vedený násilím nebo vojenskými prostředky. Kyberkriminalita řeší pachatele a trestný čin. Kyberbezpečnost řeší ochranu systémů a služeb. Kyberobrana řeší vojenskou schopnost bránit stát v kyberprostoru. Kyberválka je až situace, kdy kyberoperace plní válečné nebo vojenské cíle. Proto je přesnější říct, že současný svět je plný kyberkonfliktů, ale jen část z nich patří pod kyberválku.

Mezinárodní právo platí i v kyberprostoru. To znamená, že se uvažuje zákaz použití síly, právo na sebeobranu, mezinárodní humanitární právo, princip rozlišování, proporcionality a vojenské nezbytnosti. Sporné je, kdy přesně kyberoperace dosáhne prahu použití síly nebo ozbrojeného útoku.

\subsection*{Mezinárodní právo a kyberoperace}
U kyberválky je dobré umět říct, že nejde o právní vakuum. Státy se dlouhodobě shodují alespoň na základní tezi, že mezinárodní právo se vztahuje i na kyberprostor. Sporné jsou detaily aplikace.

\begin{itemize}
  \item \textbf{Svrchovanost} znamená, že stát má právo vykonávat autoritu nad svým územím, infrastrukturou a veřejnou mocí. Kyberoperace proti systémům na území státu může být porušením svrchovanosti, zejména pokud má vážné účinky nebo zasahuje do státních funkcí.
  \item \textbf{Zásada nevměšování} chrání rozhodovací autonomii státu. Nestačí běžná špionáž nebo propaganda; problém vzniká hlavně tehdy, když kyberoperace donucuje stát v oblasti jeho svrchovaných funkcí, například voleb, bezpečnosti nebo veřejné správy.
  \item \textbf{Zákaz použití síly} se může uplatnit, pokud kyberoperace svými účinky odpovídá použití síly, například rozsáhlé fyzické škody, ztráty na životech nebo paralyzování kritické infrastruktury. Ještě vyšším prahem je \textbf{ozbrojený útok}, který může založit právo na sebeobranu.
  \item \textbf{Retorze a protiopatření} jsou reakce pod prahem sebeobrany. Retorze je nepřátelský, ale dovolený krok, například diplomatické vyhoštění nebo sankce. Protiopatření by jinak byla protiprávní, ale mohou být dočasně dovolena jako reakce na předchozí mezinárodně protiprávní jednání jiného státu, při splnění proporcionality a dalších podmínek.
  \item \textbf{Odpovědnost státu} vyžaduje přičitatelnost jednání státu. U nestátních aktérů je klíčové, zda stát skupinu řídil, kontroloval, pověřil, toleroval způsobem právně relevantním, nebo zda později jednání převzal za své. Politická atribuce a právní přičitatelnost nejsou totéž.
  \item \textbf{Mezinárodní humanitární právo} se uplatní v ozbrojeném konfliktu. I kyberoperace pak musí respektovat rozlišování mezi civilními a vojenskými cíli, proporcionalitu, vojenskou nezbytnost a zákaz zbytečného utrpení. Civilní nemocnice, vodárny nebo energetika nejsou volným cílem jen proto, že jsou digitálně propojené.
\end{itemize}

\pravobox{Tallinn Manual je neformální expertní studie o aplikaci mezinárodního práva na kyberoperace. Není to smlouva, není právně závazný a nezakládá povinnosti států. Je ale užitečný jako mapa argumentů: ukazuje, jak lze na kyberoperace aplikovat svrchovanost, zákaz použití síly, sebeobranu, odpovědnost státu a humanitární právo.}

\pozor{NATO opakovaně uvádí, že závažný kybernetický útok může vést k uplatnění článku 5 Severoatlantické smlouvy, ale vždy case-by-case a politickým rozhodnutím spojenců. Není automatické pravidlo "kyberútok = článek 5". Zkouškově je přesné říct: může, pokud účinky a kontext dosáhnou potřebné závažnosti.}

\pozor{Silná odpověď neříká "každý útok státu je kyberválka". Řekne: záleží na aktérovi, cíli, rozsahu, účincích, politickém kontextu a napojení na ozbrojený konflikt. Špionáž může být nepřátelská a závažná, ale sama o sobě obvykle není válka.}

\subsection*{Kyberprostor, vrstvy a typy účinků}
Kyberprostor není jen "internet". Pro zkoušku je užitečné ho rozdělit na několik vrstev. \textbf{Fyzická vrstva} jsou kabely, servery, datová centra, satelity, koncová zařízení a průmyslová zařízení. \textbf{Logická/protokolová vrstva} jsou adresace, směrování, protokoly, domény, autentizace a síťová pravidla. \textbf{Datová nebo sémantická vrstva} je obsah: databáze, e-maily, dokumenty, konfigurace, účty. \textbf{Kognitivní vrstva} je lidské vnímání, rozhodování, důvěra a informační působení.

Kyberoperace mohou mít různé účinky. \textbf{Logický účinek} je změna dat, účtů, konfigurace nebo dostupnosti služby. \textbf{Fyzický účinek} nastává, když zásah do systému ovlivní reálný proces, například turbínu, tlak, signalizaci, nemocniční provoz nebo průmyslovou linku. \textbf{Kognitivní účinek} míří na důvěru, rozhodování, paniku nebo politickou legitimitu. Proto je důležitá triáda CIA - důvěrnost, integrita, dostupnost - ale u konfliktu je potřeba přidat i odolnost, obnovu a dopad na společnost.

\priklad{Když útočník ukradne diplomatické e-maily, jde hlavně o důvěrnost a politický účinek. Když změní konfiguraci průmyslového řídicího systému, jde o integritu a možný fyzický účinek. Když shodí online služby státu před volbami a současně šíří narativ o neschopnosti vlády, kombinuje dostupnost a kognitivní účinek.}

Typy kyberoperací se dají vysvětlit podle toho, co chtějí způsobit. \textbf{Akvizice informací} znamená špionáž a krádež dat. \textbf{Zveřejnění nebo šíření informací} znamená úniky, propagandu a dezinformace. \textbf{Disrupce} znamená narušení provozu, například DDoS nebo vyřazení části služby. \textbf{Destrukce} znamená mazání dat, wipery nebo výjimečně fyzické poškození zařízení. Tím se dá pěkně ukázat, že kyberválka není jen "hackování", ale práce s účinky na rozhodování, provoz a schopnost státu jednat.

\pozor{Kyberválku nepleť s elektronickým bojem. Elektronický boj pracuje s elektromagnetickým spektrem, rušením, radiolokací a signály. Může vojensky souviset s kyberoperacemi, ale není to totéž.}

\subsection*{Historie a současné trendy}
Vývoj kyberkonfliktu lze vykládat v několika milnících:
\begin{itemize}
  \item \textbf{rané incidenty} - ukázaly zranitelnost propojených systémů a význam špionáže,
  \item \textbf{Estonsko 2007} - masivní DDoS a politicky motivované útoky jako symbol zranitelnosti digitální společnosti,
  \item \textbf{Gruzie 2008} - kyberútoky doprovázející konvenční vojenský konflikt,
  \item \textbf{Stuxnet} - sabotáž průmyslových systémů, propojení kyberprostoru a fyzického světa,
  \item \textbf{Ukrajina od roku 2014} - útoky na energetiku, destruktivní malware, informační operace a testování ruských kapacit,
  \item \textbf{NotPetya 2017} - destruktivní malware maskovaný jako ransomware, globální dopady na soukromé firmy,
  \item \textbf{válka Ruska proti Ukrajině po roce 2022} - integrace kyberoperací, zpravodajství, satelitních služeb, informačních kampaní a civilní digitální infrastruktury; často se uvádí destruktivní malware HermeticWiper, útok na satelitní službu Viasat a opakované DDoS a wipery proti ukrajinským institucím.
\end{itemize}

Současné trendy jsou hlavně profesionalizace, komoditizace nástrojů, outsourcing přes proxy aktéry, ransomware jako strategická hrozba, útoky na dodavatelské řetězce, zneužívání cloudových identit, prolínání kyberoperací s dezinformacemi a využívání AI pro sociální inženýrství, automatizaci a analýzu cílů. Zároveň roste význam obrany: sdílení indikátorů, threat intelligence, zero trust, segmentace, cvičení a odolnost.

Historii nemusíš odříkat jako seznam letopočtů. Lepší je ukázat vývoj logiky konfliktu. Rané incidenty ukázaly, že propojené sítě jsou zranitelné. Estonsko ukázalo, že digitální stát může být politicky ochromen DDoS kampaní. Gruzie ukázala propojení kyberútoků s konvenčním konfliktem. Stuxnet ukázal, že kyberoperace může zasáhnout průmyslový fyzický proces. Ukrajina ukazuje dlouhodobé testování infrastruktury, wipery, satelitní komunikaci, DDoS, informační operace a propojení civilních technologií s válkou. Tím se dostaneš od "historie" k současným trendům přirozeně.

\subsection*{Aktéři}
\begin{itemize}
  \item \textbf{Státy} - mají strategické cíle, zpravodajské kapacity, peníze a dlouhodobost.
  \item \textbf{Státem podporované skupiny} - APT aktéři, kteří mohou být formálně odděleni od státu, ale plní jeho zájmy.
  \item \textbf{Proxy aktéři} - kriminální nebo hacktivistické skupiny, které stát toleruje, řídí nebo využívá.
  \item \textbf{Kyberkriminální skupiny} - motivované ziskem, ale někdy s politickým přesahem.
  \item \textbf{Hacktivisté} - motivovaní ideologií, solidaritou nebo politickým protestem.
  \item \textbf{Soukromé firmy} - poskytují bezpečnost, cloud, satelitní komunikaci, threat intelligence i ofensivní nástroje.
  \item \textbf{Jednotlivci} - menší kapacita, ale v asymetrickém prostředí mohou způsobit velký dopad.
\end{itemize}

U APT skupin se zdůrazňuje dlouhodobost, skrytost, dobré financování, cílený výběr obětí a vazba na státní zájmy. Jako příklady se v literatuře často uvádějí skupiny spojované s Ruskem, Čínou, KLDR, Íránem nebo západními službami. Není nutné memorovat desítky názvů, důležitější je vysvětlit logiku: APT neútočí jen pro rychlý zisk, ale pro špionáž, sabotáž, přípravu přístupu pro budoucí konflikt nebo politický tlak.

Hacktivismus může být autentický politický protest, ale také krytí pro státní nebo kriminální operaci. Cyberterrorismus je pojem, se kterým je lepší zacházet opatrně: mnoho útoků je politicky motivovaných, ale nedosahuje úrovně teroristického násilí. U státnic tím ukážeš, že nebereš mediální nálepky automaticky.

\subsection*{Atribuce kyberútoků}
Atribuce je přisouzení útoku konkrétnímu aktérovi. Je obtížná, protože útočníci používají proxy servery, kompromitovanou infrastrukturu, falešné stopy, sdílené nástroje a ukradený kód. Atribuce má několik rovin:
\begin{itemize}
  \item \textbf{technická} - indikátory kompromitace, malware, infrastruktura, TTPs,
  \item \textbf{operační} - jak skupina pracuje, kdy útočí, jaké cíle vybírá,
  \item \textbf{strategická} - komu útok prospívá a zapadá do politického kontextu,
  \item \textbf{právní a politická} - zda stát atribuci veřejně oznámí a s jakými důkazy.
\end{itemize}
Technicky lze mít vysokou pravděpodobnost, ale politická atribuce je rozhodnutí. Stát často nezveřejní všechny důkazy, protože by odhalil zpravodajské zdroje a metody.

Dobré vysvětlení atribuce je: "Atribuce není jen otázka, z jaké IP adresy útok přišel." IP adresa může patřit napadenému serveru třetí strany. Malware může být záměrně podobný jiné skupině. Útočník může pracovat v časech, které mají napodobit jinou zemi. Proto se skládají technické stopy, chování skupiny, cíle útoku, geopolitický kontext a zpravodajské informace. Výsledkem často není absolutní jistota, ale míra důvěry. Veřejná atribuce je navíc politický krok: stát tím vysílá signál a otevírá prostor pro sankce, diplomacii nebo kolektivní reakci.

\subsection*{Odstrašení v kyberprostoru}
Klasické jaderné odstrašení stojí na jasné atribuci, vysoké ceně útoku a věrohodné odvetě. V kyberprostoru je to složitější: atribuce je nejistá, útoky jsou pod prahem války, aktéři jsou různí a náklady na útok mohou být nízké.

Přesto odstrašení existuje:
\begin{itemize}
  \item \textbf{deterrence by punishment} - hrozba odvetou: sankce, trestní stíhání, diplomatické kroky, kybernetická protiopatření,
  \item \textbf{deterrence by denial} - ztížení útoku tak, aby se nevyplatil: segmentace, MFA, zálohy, monitoring,
  \item \textbf{resilience} - schopnost rychle obnovit provoz a snížit strategický efekt útoku,
  \item \textbf{normy a reputace} - budování pravidel odpovědného chování států,
  \item \textbf{kolektivní reakce} - NATO, EU, společné atribuce, sankční režimy.
\end{itemize}
V praxi je nejúčinnější kombinace. Samotná hrozba odvetou nestačí, pokud útočník věří, že nebude odhalen nebo že oběť nebude reagovat.

\pozor{U odstrašení vždy zmiň praktické otázky: Lze útok spolehlivě přisoudit? Je odveta věrohodná? Nepovede reakce k eskalaci? Má oběť kapacity reagovat? Nevyužil útočník infrastrukturu třetího státu? A není levnější útok prostě přežít díky dobré odolnosti?}

Na příkladu nemocnice nebo energetiky je rozdíl mezi typy odstrašení vidět jednoduše. Odstrašení trestem říká: když zaútočíš, vystavuješ se sankcím, trestnímu stíhání, diplomatické reakci nebo protiopatřením. Odstrašení odepřením říká: útok se ti nevyplatí, protože máme segmentaci, zálohy, monitoring a obnovu. Odolnost říká: i když něco zasáhneš, nezískáš strategický efekt, protože službu obnovíme a společnost se nezhroutí. V kyberprostoru je často nejrealističtější právě kombinace odepření a odolnosti.

\subsection*{Netwars}
Koncept netwars spojovaný hlavně s Johnem Arquillou a Davidem Ronfeldtem popisuje konflikty vedené síťově organizovanými aktéry, kteří využívají informační technologie, decentralizaci, rychlou komunikaci a propojení mnoha uzlů. Nejde jen o hacking. Netwar může zahrnovat propagandu, mobilizaci, dezinformace, psychologické působení, hacktivismus, úniky dat, koordinované kampaně a tlak na veřejné mínění.

Rozdíl je tento: \textbf{kyberválka} je užší pojem pro kybernetické prostředky v ozbrojeném nebo vojensky významném konfliktu. \textbf{Netwar} je širší síťový konflikt o informace, legitimitu, mobilizaci a rozhodování. Netwar může kyberválku doprovázet, ale může probíhat i bez ní, například jako dlouhodobá dezinformační, hacktivistická nebo úniková kampaň.

V kyberkonfliktu se netwars projevují tak, že vedle technických útoků běží informační operace. Útočník nejen shodí službu, ale zároveň šíří narativ, že stát je neschopný. Nejen ukradne e-maily, ale zveřejní je ve správný politický moment. Nejen napadne infrastrukturu, ale posiluje chaos na sociálních sítích.

Síťové konflikty mohou mít různé organizační podoby. \textbf{Řetězová síť} předává úkoly po linii, \textbf{hvězdicová síť} má centrální uzel, \textbf{all-channel síť} umožňuje komunikaci mezi všemi uzly. Právě poslední model je typický pro decentralizované kampaně, kde aktéři sdílejí nástroje, cíle, narativy a výsledky. Častá je taktika \textbf{swarmingu}: mnoho menších aktérů působí současně na více cílů, takže obránce zahlcuje nejen technický provoz, ale i komunikace, reputační krize a rozhodování.

Netwar je dobré vysvětlit jako konflikt sítí proti hierarchii. Tradiční stát rozhoduje pomaleji, přes instituce, procedury a odpovědnosti. Síťově organizovaní aktéři mohou jednat rychle, decentralizovaně a často bez jasného velitele. To je výhoda u kampaní, kde se kombinuje únik dat, DDoS, memy, dezinformace, dobrovolnické skupiny, kriminální infrastruktura a mediální tlak. Obrana pak není jen technická, ale i komunikační a politická: stát musí chránit systémy, vysvětlovat situaci veřejnosti a nenechat si vnutit útočníkův narativ.

\priklad{Před volbami dojde k průniku do e-mailů politické strany, selektivnímu úniku dokumentů, DDoS na volební web a koordinované kampani na sociálních sítích. Technická část je kyberoperace, celek je síťový informační konflikt - netwar.}

\zkouska{V odpovědi zdůrazni hranici: kyberválka je úzký pojem, kyberkonflikt širší. Nejčastější realita je působení pod prahem války, v šedé zóně, s nejistou atribucí a kombinací technických a informačních prostředků.}

\zdroje{Přiložené PDF a BSS poznámky; Tallinn Manual jako akademický rámec; NATO a EU dokumenty k odpovědnému chování a kolektivní reakci.}

\newpage

\otazka{18. Ochrana kritické infrastruktury}{Kyberútoky na kritickou infrastrukturu. Přisouzení a odstrašování v kyberprostoru.}{BSSb1103}

\subsection*{Kostra odpovědi}
\begin{enumerate}
  \item Vymezit kritickou infrastrukturu a kritické subjekty: služby nezbytné pro fungování státu a společnosti.
  \item Dát aktuální právní kontext: CER a zákon o odolnosti subjektů kritické infrastruktury řeší odolnost kritických subjektů; zákon o kybernetické bezpečnosti řeší kybernetickou bezpečnost regulovaných služeb.
  \item Vysvětlit sektory, interdependence a rozdíl IT/OT/ICS/SCADA.
  \item Popsat typické kyberútoky na KI a konkrétní příklady.
  \item Vyložit ochranu: prevence, detekce, reakce, obnova, kontinuita, fyzická bezpečnost, dodavatelé a cvičení.
  \item Zopakovat atribuci a odstrašení v kontextu KI, včetně prahu reakce a rizika eskalace.
\end{enumerate}

\subsection*{Vymezení kritické infrastruktury}
Kritická infrastruktura zahrnuje prvky, systémy a služby, jejichž narušení by mělo závažný dopad na bezpečnost státu, životy a zdraví osob, ekonomiku nebo základní funkce společnosti. Typicky jde o energetiku, dopravu, zdravotnictví, vodní hospodářství, digitální infrastrukturu, finanční služby, veřejnou správu, potraviny, chemický průmysl a další základní odvětví.

Základní myšlenka je jednoduchá: nejde jen o to, že existuje důležitá budova, kabel nebo server. Důležité je, jaká \emph{základní služba} na nich závisí. Výpadek nemocničního informačního systému není problém proto, že "nejde počítač", ale proto, že se zpozdí péče, nejdou plánovat operace, laboratoře neposílají výsledky a pacienti se mohou dostat do reálného ohrožení.

\pravobox{Aktuálně je potřeba rozlišovat dvě vrstvy. \textbf{CER} - směrnice 2022/2557 - je v ČR transponována zákonem o odolnosti subjektů kritické infrastruktury (č. 266/2025 Sb.), účinným od 19. 8. 2025. Ten mění důraz ze starého modelu "stát určí prvky kritické infrastruktury" na model "určí se subjekty poskytující základní služby a ty samy určují svou kritickou infrastrukturu, rizika, dodavatele a plán odolnosti". \textbf{ZKB/NIS2} - zákon o kybernetické bezpečnosti (č. 264/2025 Sb.) a vyhlášky - řeší kybernetickou rovinu regulovaných služeb, vyšší/nižší režim, bezpečnostní opatření a hlášení incidentů. Tyto režimy se překrývají, ale nejsou totožné: CER je all-hazards odolnost, ZKB je kybernetická bezpečnost.}

\subsection*{Režim CER/ZKI v praxi}
Nová úprava kritické infrastruktury se dá vysvětlit jako řetězec:
\begin{enumerate}
  \item \textbf{Stát posoudí rizika} - nejprve se na národní úrovni hodnotí, které základní služby jsou pro stát a společnost kritické a jaké hrozby je mohou narušit.
  \item \textbf{Stát určí kritické subjekty} - podle odvětví, základní služby a významových kritérií se rozhoduje, kdo spadá do režimu kritického subjektu.
  \item \textbf{Subjekt posoudí vlastní rizika} - kritický subjekt si musí vyhodnotit, co jeho základní službu ohrožuje: kyberútok, fyzická sabotáž, výpadek energie, dodavatel, povodeň, personální nedostatek nebo kombinace.
  \item \textbf{Subjekt zavede opatření} - organizační, technická, fyzická a personální opatření k prevenci, detekci, reakci a obnově.
  \item \textbf{Subjekt připraví plán odolnosti} - popíše kritickou infrastrukturu, závislosti, odpovědnosti, krizové scénáře, komunikaci a obnovu služby.
  \item \textbf{Incident se řeší a hlásí} - subjekt musí umět incident zvládnout, dokumentovat a komunikovat s příslušnými orgány.
  \item \textbf{Stát dohlíží a koordinuje} - kontroluje plnění povinností, sdílí metodiku a propojuje sektorové gestory.
\end{enumerate}

Klíčové slovo je \textbf{all-hazards}. CER/ZKI neřeší jen hackery. Stejný plán odolnosti má dávat smysl pro ransomware, požár, povodeň, blackout, sabotáž, pandemii, selhání dodavatele i kombinovaný hybridní scénář. Proto se liší od čistě kybernetické regulace: ZKB se ptá na bezpečnost informačních a komunikačních systémů regulované služby, zatímco ZKI se ptá, zda základní služba přežije široké spektrum narušení.

\pozor{K 16. 6. 2026 je potřeba mluvit opatrně o přechodném období. Zákon o odolnosti subjektů kritické infrastruktury je účinný od 19. 8. 2025, ale praktické určování kritických subjektů je navázané na prováděcí úpravu základních služeb a významových kritérií. Staré prvky kritické infrastruktury proto nelze jednoduše prohlásit za okamžitě nahrazené plně fungujícím novým systémem. Před zkouškou ověř, zda už je účinné příslušné nařízení vlády, jaké běží lhůty a které subjekty byly nově určeny.}

Starší poznámky často pracují s pojmy \textbf{kritická komunikační infrastruktura} a \textbf{kritická informační infrastruktura}. U zkoušky je můžeš použít jako vysvětlující zkratky, ale s aktuální korekcí. Komunikační infrastruktura znamená sítě a služby elektronických komunikací, které umožňují přenos informací: pevné a mobilní sítě, páteřní internet, datová centra, satelitní komunikace nebo veřejné radiokomunikační sítě. Informační infrastruktura znamená informační systémy a data, bez nichž nemůže fungovat veřejná správa, zdravotnictví, energetika, finance nebo doprava. Nový ZKB už odpověď staví hlavně přes regulovanou službu a jejího poskytovatele, ale věcně pořád chráníme totéž: systémy, sítě, data a služby, jejichž výpadek má společenský dopad.

\subsection*{Sektory a vzájemné závislosti}
Kritická infrastruktura není soubor izolovaných ostrovů. Sektory jsou propojené a právě to vytváří systémové riziko:
\begin{itemize}
  \item \textbf{energetika} podpírá téměř všechno ostatní: datová centra, nemocnice, vodárenství, dopravu i bankovnictví,
  \item \textbf{digitální a komunikační infrastruktura} umožňuje řízení sítí, komunikaci složek státu, platby, logistiku a krizové řízení,
  \item \textbf{zdravotnictví} závisí na elektřině, datech, přístrojích, dodavatelích léčiv a schopnosti komunikovat,
  \item \textbf{finance} závisí na konektivitě, identitě, dostupnosti platebních systémů a důvěře,
  \item \textbf{doprava a logistika} závisí na plánovacích systémech, GPS, přístavech, letištích, železnici, palivech a komunikaci,
  \item \textbf{vodohospodářství a potraviny} jsou méně viditelné, ale jejich výpadek rychle dopadá na základní životní potřeby.
\end{itemize}
Proto se u KI mluví o kaskádových dopadech. Útok na jednu službu může způsobit řetězovou reakci v jiných sektorech. Výpadek elektřiny poškodí telekomunikace; výpadek telekomunikací komplikuje obnovu elektřiny; výpadek plateb nebo logistiky zasáhne zásobování.

\subsection*{Proč je kyberútok na KI zvláštní}
Útok na běžnou firmu může znamenat finanční škodu. Útok na kritickou infrastrukturu může zastavit péči, dodávky elektřiny, dopravu, komunikaci nebo fungování úřadů. U KI se proto neřeší jen důvěrnost dat, ale hlavně dostupnost, integrita provozu a bezpečnost lidí.

Specifické jsou průmyslové systémy OT/ICS/SCADA. \textbf{IT} je klasické informační prostředí: servery, stanice, aplikace, databáze. \textbf{OT} je provozní technologie: řídí fyzický proces, například tlak, průtok, spínače, ventily, turbíny nebo signalizaci. \textbf{ICS} jsou průmyslové řídicí systémy a \textbf{SCADA} je dohledové a řídicí prostředí pro distribuované provozy.
\begin{itemize}
  \item mají dlouhou životnost a často běží na starších technologiích,
  \item nelze je snadno vypnout kvůli aktualizaci,
  \item výpadek může způsobit fyzickou škodu,
  \item priorita je bezpečný provoz, ne jen IT čistota,
  \item bývají závislé na dodavatelích a vzdálené správě.
\end{itemize}
Proto se ochrana OT nedá kopírovat z kancelářského IT. V běžném IT je často prioritou důvěrnost a integrita; v OT je často první dostupnost a bezpečný fyzický provoz. Potřebuje segmentaci, řízení vzdáleného přístupu, monitoring průmyslových protokolů, testované manuální režimy, bezpečnost změn a silnou spolupráci provozu s bezpečností.

\subsection*{Typické útoky}
\begin{itemize}
  \item \textbf{Ransomware} - zašifruje IT a nepřímo zastaví provoz; u nemocnic ohrožuje péči.
  \item \textbf{Destruktivní malware} - cílem není výkupné, ale zničení nebo chaos.
  \item \textbf{Útok na OT/ICS} - manipulace s řídicími systémy, senzory nebo logikou provozu.
  \item \textbf{DDoS} - nedostupnost služeb veřejnosti nebo podpůrných systémů.
  \item \textbf{Dodavatelský řetězec} - kompromitace aktualizace, dodavatele nebo vzdálené správy.
  \item \textbf{Špionáž} - dlouhodobé mapování infrastruktury pro budoucí krizi.
\end{itemize}

\subsection*{Příklady, které se dobře říkají u zkoušky}
\begin{itemize}
  \item \textbf{Stuxnet / Natanz 2010} - škodlivý kód zasáhl průmyslové řídicí systémy íránského jaderného programu. Je důležitý jako příklad prolomení "kybernetické bariéry": digitální operace způsobila fyzický efekt na centrifugách a vyžadovala znalost konkrétní technologie.
  \item \textbf{Ukrajinská energetika 2015} - útok spojený s malwarem BlackEnergy vedl k výpadkům elektřiny pro statisíce odběratelů. Důležité je, že útočníci neútočili jen na techniku, ale i na schopnost reakce, například zatížením call center.
  \item \textbf{Colonial Pipeline 2021} - ransomware proti podnikové IT části vedl k preventivnímu zastavení provozu palivového potrubí. Ukazuje, že i když není přímo kompromitované OT, rozhodnutí z bezpečnostních a provozních důvodů může zastavit fyzickou službu.
  \item \textbf{České nemocnice 2020} - útoky na zdravotnická zařízení během pandemie ukázaly, že zdravotnictví je bezpečnostní téma, ne jen lokální IT problém. Ransomware ve zdravotnictví má přímý dopad na lidskou bezpečnost.
  \item \textbf{Bangladesh Bank 2016} - útok na finanční infrastrukturu a SWIFT ukazuje, že KI není jen elektrárna nebo nemocnice; finanční systém je kritický kvůli důvěře, platbám a přeshraničním vazbám.
\end{itemize}

\priklad{Útok na energetickou distribuční společnost nemusí okamžitě vypnout elektřinu. Může nejdřív mapovat síť, ukrást přístupy, narušit dohledové systémy a připravit možnost škody v době politické nebo vojenské krize. U KI proto není špionáž "méně důležitá" jen proto, že zatím nic nespadlo; může jít o přípravu budoucí sabotáže.}

\subsection*{Ochrana kritické infrastruktury}
Ochrana KI je kombinace prevence, detekce, reakce a obnovy. Nestačí koupit bezpečnostní nástroj; cílem je odolnost služby.
\begin{itemize}
  \item \textbf{řízení rizik} - vědět, které služby jsou kritické a na čem závisejí,
  \item \textbf{mapování závislostí} - znát dodavatele, datové toky, vzdálené přístupy, záložní trasy a single points of failure,
  \item \textbf{segmentace} - oddělit kancelářské IT, kritické systémy a OT,
  \item \textbf{přístupová práva} - MFA, privilegované účty, princip nejmenších oprávnění,
  \item \textbf{zálohy a obnova} - offline/immutable zálohy a pravidelné testy obnovy,
  \item \textbf{monitoring a detekce} - SIEM, EDR, logy, OT monitoring,
  \item \textbf{dodavatelská bezpečnost} - smlouvy, kontrola vzdálených přístupů, bezpečnostní požadavky,
  \item \textbf{BCM a krizové plány} - kontinuita služeb, manuální postupy, cvičení,
  \item \textbf{fyzická bezpečnost} - přístup do objektů, napájení, ochrana datových center.
  \item \textbf{cvičení} - table-top cvičení, technická cvičení, red teaming a kybernetické polygony; cvičí se rozhodování, komunikace i technická reakce.
\end{itemize}

CER/ZKI do toho přidává jazyk odolnosti: posouzení rizik, plán odolnosti, určení manažera kritické infrastruktury, určení vlastní kritické infrastruktury a kritických dodavatelů, hlášení incidentů, cvičení a komunikaci s věcně příslušnými gestory. ZKB přidává kybernetická opatření a incident reporting. Dohromady to znamená, že organizace musí umět odpovědět na otázku: \emph{jakou základní službu poskytujeme, co ji může zastavit, jak tomu bráníme a jak ji obnovíme, když obrana selže?}

\subsection*{Přisouzení a odstrašování u KI}
Atribuce u útoků na KI je politicky citlivá, protože veřejné přisouzení státu může znamenat eskalaci. Technické stopy samy o sobě často nestačí: útočník může použít VPN, proxy, botnet, kompromitovanou infrastrukturu, falešné stopy, ukradený malware nebo časové značky napodobující jinou skupinu. Proto se atribuce skládá z techniky, zpravodajských informací, geopolitického kontextu, motivu a posouzení, komu útok prospívá.

U KI je problém ještě ostřejší: před odvetou musí stát vědět, zda jde o kriminální ransomware, státem tolerovanou skupinu, proxy aktéra, vojenskou operaci nebo kombinaci. Odveta proti špatnému cíli neodstraší, ale vytvoří další konflikt. I když se útočník technicky přisoudí, nemusí být veřejně zveřejněny všechny důkazy, protože by se odhalily zpravodajské zdroje a metody.

Odstrašení v kyberprostoru má několik forem:
\begin{itemize}
  \item \textbf{odstrašení trestem} - sankce, trestní stíhání, diplomatická reakce, veřejná atribuce, případně kybernetická nebo jiná protiopatření,
  \item \textbf{odstrašení znemožněním} - útok se nevyplatí, protože organizace má segmentaci, zálohy, monitoring, obnovení provozu a silnou reakci,
  \item \textbf{resilience} - útočník sice zasáhne, ale nedosáhne strategického efektu, protože služba pokračuje nebo se rychle obnoví,
  \item \textbf{kolektivní reakce} - EU/NATO, společné atribuce, koordinované sankce a sdílení informací,
  \item \textbf{normy a právo} - budování hranic přijatelného chování, zejména u útoků na civilní a zdravotnickou infrastrukturu.
\end{itemize}

Starší debata o odstrašení klade dobré otázky: Kde je práh reakce? Je to počet obětí, rozsah výpadku, finanční škoda, zásah do zdravotnictví nebo útok během ozbrojeného konfliktu? Dokáže stát reagovat opakovaně, nebo jednorázová odveta odhalí vlastní schopnosti? Nezapojí se třetí strany, které využijí odhalenou zranitelnost? Přijme veřejnost sílu odvety, nebo ji bude považovat za eskalaci? U KI proto dává největší smysl kombinace: dobrá obrana, rychlá obnova, věrohodná atribuce, promyšlená komunikace a škála reakcí od trestního řízení po mezinárodní sankce.

\pozor{Kyberterorismus u KI zmiň opatrně. Terorismus potřebuje politický cíl, zastrašení, publicitu a symbolický efekt. Mnoho útoků na infrastrukturu je závažných, ale právně a analyticky půjde spíš o kriminalitu, špionáž, sabotáž nebo státní/proxy operaci než o terorismus.}

\zkouska{U KI vždy řekni "nejde jen o data, ale o kontinuitu základní služby a fyzické dopady". Potom přidej právní rozlišení: CER/ZKI = odolnost kritických subjektů vůči všem hrozbám, ZKB/NIS2 = kybernetická bezpečnost regulovaných služeb.}

\zdroje{Přiložené PDF; starší výpisky k otázce 4; směrnice CER 2022/2557; zákon o odolnosti subjektů kritické infrastruktury (č. 266/2025 Sb.); zákon o kybernetické bezpečnosti (č. 264/2025 Sb.); HZS ČR k nové úpravě kritické infrastruktury; NÚKIB metodiky k regulovaným službám.}

\newpage

\otazka{19. Právní úprava kyberbezpečnosti v ČR a EU}{Základní instituty, principy, povinné orgány, systém zajištění kybernetické bezpečnosti.}{BVV03K}

\subsection*{Kostra odpovědi}
\begin{enumerate}
  \item Odlišit kybernetickou bezpečnost, kyberobranu, kyberkriminalitu a ochranu osobních údajů.
  \item Vyložit EU rámec: NIS2, CER, DORA, CRA a související předpisy.
  \item Vyložit český rámec: zákon o kybernetické bezpečnosti a prováděcí vyhlášky.
  \item Popsat regulované služby, vyšší a nižší režim, sebeidentifikaci a hlášení incidentů.
  \item Popsat orgány a systém zajištění kybernetické bezpečnosti.
\end{enumerate}

\subsection*{Základní rozlišení}
Kybernetická bezpečnost je prevence, řízení rizik, ochrana systémů a služeb, detekce incidentů, reakce a obnova. Kyberobrana je vojenská dimenze kyberprostoru a patří do obranné politiky. Kyberkriminalita je trestná činnost a řeší ji orgány činné v trestním řízení. Ochrana osobních údajů chrání práva subjektů údajů a řeší ji hlavně GDPR a ÚOOÚ.

Jeden incident může spadat do více režimů. Ransomware v nemocnici může být kybernetický bezpečnostní incident podle ZKB, trestný čin podle trestního zákoníku, porušení zabezpečení osobních údajů podle GDPR a krizová situace z hlediska zdravotnictví.

\subsection*{EU rámec}
\textbf{NIS2} je hlavní směrnice EU pro vysokou společnou úroveň kybernetické bezpečnosti. Rozšiřuje okruh regulovaných subjektů, posiluje řízení rizik, odpovědnost vedení, hlášení incidentů, dohled a sankce. V ČR je transponována novým zákonem o kybernetické bezpečnosti.

\textbf{CER} řeší odolnost kritických subjektů v širším smyslu: fyzická, organizační a provozní odolnost. Je důležitá tam, kde narušení služby ohrožuje společnost bez ohledu na to, zda příčinou je kyberútok, sabotáž, přírodní událost nebo technické selhání.

\textbf{DORA} je zvláštní úprava pro finanční sektor. Zaměřuje se na řízení ICT rizik, hlášení incidentů, testování digitální provozní odolnosti a rizika spojená s externími ICT poskytovateli. U zkoušky je dobré říct, že DORA je lex specialis pro finance, ale neznamená, že finanční subjekt nikdy neřeší jiné předpisy.

\textbf{CRA} zavádí požadavky na produkty s digitálními prvky. Posouvá kyberbezpečnost do fáze návrhu, vývoje a uvádění produktů na trh. Výrobci budou řešit bezpečnostní požadavky, správu zranitelností, aktualizace a hlášení aktivně zneužívaných zranitelností.

Vedle těchto předpisů je dobré zmínit i institucionální vrstvu EU. \textbf{ENISA} podporuje členské státy, certifikaci, metodiky a evropskou spolupráci. \textbf{CSIRT Network} propojuje národní CSIRT týmy. \textbf{EU-CyCLONe} podporuje koordinaci rozsáhlých kybernetických krizí. \textbf{CERT-EU} chrání instituce EU. U trestních věcí vstupuje \textbf{Europol/EC3} a Eurojust. EU tedy nedělá jen normy, ale buduje i praktické koordinační a krizové mechanismy.

\textbf{Cybersecurity Act} - nařízení 2019/881 - je důležitý proto, že dal ENISA stálý mandát jako Agentuře EU pro kybernetickou bezpečnost a vytvořil evropský rámec kyberbezpečnostní certifikace. Certifikace se týká ICT produktů, služeb a procesů; může být dobrovolná, ale jiné předpisy nebo zadávací podmínky z ní mohou udělat prakticky nutný standard. Zkouškově si pamatuj: NIS2 řeší povinnosti subjektů, Cybersecurity Act řeší ENISA a certifikační rámec.

\textbf{Cyber Solidarity Act} - nařízení 2025/38 - posiluje unijní schopnosti detekce, připravenosti a reakce na velké kybernetické hrozby a incidenty. Staví na třech pilířích: evropský systém kybernetického varování včetně přeshraničních bezpečnostních operačních center, mechanismus kybernetické nouze včetně evropské kybernetické rezervy a mechanismus přezkumu incidentů. Nejde o náhradu NIS2, ale o operační a krizovou nadstavbu pro společnou detekci, podporu a poučení z rozsáhlých incidentů.

\subsection*{Český zákon o kybernetické bezpečnosti}
Zákon o kybernetické bezpečnosti (č. 264/2025 Sb.) je nový základ české kyberbezpečnostní regulace. Upravuje práva a povinnosti osob a orgánů veřejné moci v oblasti zajišťování kybernetické bezpečnosti, působnost NÚKIB a systém spolupráce.

Prováděcí rámec tvoří zejména \textbf{vyhláška č. 408/2025 Sb. o regulovaných službách}, \textbf{vyhláška č. 409/2025 Sb. o bezpečnostních opatřeních pro vyšší režim} a \textbf{vyhláška č. 410/2025 Sb. o bezpečnostních opatřeních pro nižší režim}. Prakticky je dobré znát jejich funkci, ne jen čísla: první pomáhá určit, kdo spadá do regulace a v jakém režimu; druhá a třetí říkají, jaká organizační a technická opatření se mají zavést. Pro veřejnou správu a cloud jsou navíc relevantní předpisy k bezpečnostním úrovním informačních systémů veřejné správy a bezpečnostním pravidlům pro cloud computing.

Praktické jádro je v určení regulovaných služeb a režimů:
\begin{itemize}
  \item subjekt musí posoudit, zda poskytuje regulovanou službu v regulovaném odvětví,
  \item pokud ano, řeší se velikost, význam služby a další kritéria podle vyhlášky o regulovaných službách,
  \item výsledek vede k režimu vyšších nebo nižších povinností,
  \item subjekt provádí sebeidentifikaci a ohlašuje regulovanou službu NÚKIB,
  \item po zápisu plní organizační a technická opatření, hlásí incidenty a udržuje kontaktní údaje.
\end{itemize}

Vyšší režim je přísnější: detailnější bezpečnostní opatření, audity a vyšší nároky na řízení. Nižší režim je jednodušší, ale stále vyžaduje základní řízení bezpečnosti, odpovědnosti, opatření a hlášení.

Nová česká úprava je postavená na principu sebeidentifikace. Subjekt si nemá jen pasivně čekat na dopis od státu, ale musí vyhodnotit, zda poskytuje regulovanou službu a zda spadá do vyššího nebo nižšího režimu. Prakticky se řeší zejména odvětví, velikost subjektu, význam poskytované služby, dopad jejího výpadku a zvláštní kritéria podle prováděcí vyhlášky. Po ohlášení a zápisu se rozbíhá sada povinností: kontaktní údaje, bezpečnostní role, řízení rizik, dokumentace, školení, řízení dodavatelů, evidence aktiv, technická opatření a hlášení incidentů.

\subsection*{Strategicky významné služby a dodavatelský řetězec}
Nový ZKB řeší i bezpečnost dodavatelského řetězce. Neznamená to, že stát plošně prověřuje každého dodavatele každé regulované organizace. Speciální mechanismus bezpečnosti dodavatelského řetězce se váže hlavně na \textbf{strategicky významné služby} a na dodavatele, kteří mohou mít zásadní vliv na jejich bezpečnost. Smyslem je řešit nejen technickou zranitelnost produktu, ale i strategické riziko: vazby dodavatele na rizikový stát, možnost nátlaku, právní prostředí, vlastnickou strukturu, důvěryhodnost a schopnost dlouhodobé podpory.

Stát přitom nemá jednat libovolně. Opatření musí být opřená o posouzení rizik, zákonný postup, proporcionalitu a možnost přezkumu. V praxi může jít o omezení nebo zákaz použití určitého dodavatele u strategicky významné služby, ale až po posouzení konkrétního rizika a dopadu. U zkoušky je dobré říct: dodavatelský řetězec není jen "software má chybu"; je to kombinace technického, právního, ekonomického a geopolitického rizika.

\subsection*{Bezpečnostní opatření a nástroje NÚKIB}
Bezpečnostní opatření se tradičně dělí na \textbf{organizační} a \textbf{technická}. Organizační opatření jsou například bezpečnostní politika, řízení rizik, odpovědnosti, klasifikace aktiv, řízení dodavatelů, plán kontinuity, školení, audit a dokumentace. Technická opatření jsou například řízení přístupů, vícefaktorové ověřování, logování, monitoring, zálohování, segmentace, šifrování, ochrana koncových zařízení, detekce incidentů a testování obnovy.

NÚKIB má vedle běžného dohledu i mimořádné nástroje. \textbf{Varování} upozorňuje na hrozbu nebo zranitelnost a regulované subjekty ho musí zohlednit v řízení rizik. \textbf{Opatření} a \textbf{protiopatření} mohou ukládat konkrétní kroky ke snížení rizika nebo zvládnutí incidentu. \textbf{Stav kybernetického nebezpečí} je mimořádný režim pro závažné ohrožení kybernetické bezpečnosti; podle aktuální úpravy ho vyhlašuje ředitel NÚKIB nejdéle na 30 dnů a lze ho prodloužit tak, aby celková doba nepřekročila 60 dnů.

\pravobox{Principy regulace: \textbf{technologická neutralita} znamená, že právo nemá být navázané na jednu konkrétní technologii. \textbf{Proporcionalita} znamená, že povinnosti a zásahy musí odpovídat riziku a dopadu. \textbf{Risk-based approach} znamená, že organizace má řídit skutečná rizika, ne jen formálně plnit seznam. \textbf{Odpovědnost vedení} znamená, že kyberbezpečnost není jen problém IT oddělení. \textbf{Autonomie regulovaného subjektu} znamená, že stát stanoví rámec, ale subjekt si musí prakticky nastavit vlastní řízení bezpečnosti.}

\subsection*{Hlášení incidentů}
NIS2 logika stojí na rychlém hlášení významných incidentů, ale český ZKB rozlišuje vyšší a nižší režim. \textbf{Vyšší režim} je přísnější: poskytovatel hlásí kybernetické bezpečnostní incidenty, které postihují regulovanou službu, mají původ v kyberprostoru a nelze u nich vyloučit úmyslné zavinění; prvotní hlášení má být rychlé, typicky do 24 hodin. NÚKIB následně posuzuje, zda jde o incident s významným dopadem. \textbf{Nižší režim} je užší: poskytovatel typicky hlásí až incidenty s významným dopadem, který vyhodnocuje podle vyhlášky.

Pokud incident dosáhne významného dopadu, navazuje další rytmus: prvotní informace, podrobnější hlášení do 72 hodin, případné průběžné aktualizace na žádost a závěrečná zpráva typicky do 30 dnů, případně později u dlouhotrvajícího incidentu. Přesný postup se řídí zákonem a prováděcími předpisy. Zkouškově stačí umět vysvětlit logiku: stát potřebuje rychle vědět, že se něco děje, potom potřebuje doplnit technický a dopadový obraz a nakonec chce závěrečnou analýzu a poučení.

Hlášení incidentu není totéž co oznámení porušení zabezpečení osobních údajů podle GDPR. Pokud incident obsahuje únik osobních údajů, běží vedle kybernetického hlášení i GDPR režim: dozorovému úřadu zpravidla do 72 hodin od zjištění, subjektům údajů při vysokém riziku.

\subsection*{Orgány a systém}
\begin{itemize}
  \item \textbf{NÚKIB} - gestor kybernetické bezpečnosti, dohled, varování, opatření, metodiky, GovCERT.CZ.
  \item \textbf{GovCERT.CZ} - vládní tým pro řešení incidentů.
  \item \textbf{CSIRT.CZ} - národní CSIRT, koordinace incidentů v širším ekosystému.
  \item \textbf{Regulované subjekty} - nesou primární odpovědnost za bezpečnost svých služeb.
  \item \textbf{Vrcholové vedení} - schvaluje a nese odpovědnost za bezpečnostní opatření; NIS2 posiluje osobní odpovědnost managementu.
  \item \textbf{Policie, zpravodajské služby, armáda, ÚOOÚ, ČTÚ, DIA, NBÚ a další} - vstupují podle povahy incidentu a právního režimu.
\end{itemize}

Prakticky se dá systém popsat jako vícevrstvý. NÚKIB a GovCERT řeší civilní kyberbezpečnost a regulaci. CSIRT.CZ podporuje širší technickou komunitu. Policie a státní zastupitelství řeší trestní rovinu. Zpravodajské služby pomáhají s hrozbami, špionáží a atribucí. Armáda a Vojenské zpravodajství řeší obrannou dimenzi. ÚOOÚ vstupuje při osobních údajích, ČTÚ u elektronických komunikací, DIA u digitálních služeb státu a NBÚ u utajovaných informací.

\priklad{Výrobní podnik nejdřív neřeší "jsme NIS2?", ale konkrétně: poskytujeme regulovanou službu podle české vyhlášky? V jakém odvětví? Splňujeme velikostní nebo významová kritéria? Až potom určí režim povinností.}

\priklad{Nemocnice ve vyšším režimu musí uvažovat nejen antivir a firewall, ale i odpovědnost vedení, řízení dodavatelů nemocničního systému, zálohy oddělené od produkce, krizové postupy pro výpadek IT, školení personálu, hlášení incidentu NÚKIB a souběžně GDPR režim, pokud uniknou údaje pacientů.}

\pozor{Neříkej, že NIS2 přímo ukládá české firmě všechny povinnosti. Jako směrnice potřebuje transpozici; v ČR je klíčový zákon o kybernetické bezpečnosti (č. 264/2025 Sb.) a prováděcí vyhlášky.}

\zdroje{Směrnice 2022/2555 (NIS2); zákon o kybernetické bezpečnosti (č. 264/2025 Sb.); vyhlášky k regulovaným službám, vyšším povinnostem a nižším povinnostem; DORA 2022/2554; CRA 2024/2847.}

\newpage

\otazka{20. Kyberkriminalita}{Prameny práva (národní, evropské a mezinárodní), typická trestná činnost, klasifikace trestných činů, právní kvalifikace a související postupy a kritéria, mezinárodní spolupráce.}{BVV03K}

\subsection*{Kostra odpovědi}
\begin{enumerate}
  \item Vymezit kyberkriminalitu a odlišit ji od kyberbezpečnosti a kyberobrany.
  \item Vysvětlit tři významy "klasifikace": právní třídění podle skutkových podstat, technická/incidentní taxonomie a právní kvalifikace konkrétního skutku.
  \item Popsat prameny práva: národní, evropské a mezinárodní.
  \item Rozebrat typickou trestnou činnost: cyber-dependent, cyber-enabled a podpůrně cyber-supported.
  \item Ukázat právní kvalifikaci: skutková podstata, objekt, objektivní stránka, subjekt, subjektivní stránka, škoda, souběh, stadia.
  \item Doplnit související postupy: detekce, prověřování, elektronické místo činu, fyzická prohlídka, důkazy, znalec, soud.
  \item Popsat mezinárodní spolupráci: MLA, EVP, Budapešť 24/7, Europol/EC3, Eurojust, EPOC/EPOC-PR.
\end{enumerate}

\subsection*{Vymezení a vztah ke kyberbezpečnosti}
Kyberkriminalita je trestná činnost páchaná v kyberprostoru, prostřednictvím informačních a komunikačních technologií nebo přímo proti počítačovým systémům, sítím a datům. Má ji na starosti hlavně trestní justice: Policie ČR, státní zastupitelství a soudy. U závažné a organizované kriminality vstupuje Národní centrála proti organizovanému zločinu, u technických úkonů specialisté, znalci a poskytovatelé služeb.

Je dobré hned odlišit tři blízké pojmy. \textbf{Kyberbezpečnost} chrání systémy, data a služby a řeší ji hlavně NÚKIB, regulované subjekty, CERT/CSIRT týmy a správci systémů. \textbf{Kyberobrana} je obranná a vojenská dimenze, typicky Ministerstvo obrany, Vojenské zpravodajství, NCKO a kybernetické síly. \textbf{Kyberkriminalita} řeší trestní odpovědnost pachatele a procesní zajištění důkazů.

\priklad{Ransomware v nemocnici může být současně trestný čin pachatele, kybernetický bezpečnostní incident nemocnice podle ZKB, porušení zabezpečení osobních údajů podle GDPR a krizový problém zdravotnictví. U zkoušky proto vždy řekni, který režim právě řešíš.}

\subsection*{Tři významy klasifikace}
Jednotná klasifikace kyberkriminality neexistuje, takže je dobré ukázat, že slovo "klasifikace" může znamenat tři různé věci:
\begin{itemize}
  \item \textbf{právní klasifikace} - třídění podle skutkových podstat trestního zákoníku; je méně technicky detailní, ale rozhoduje o trestnosti,
  \item \textbf{technická nebo incidentní taxonomie} - popisuje typ útoku nebo incidentu: malware, phishing, DDoS, sběr informací, zranitelnost, únik dat,
  \item \textbf{právní kvalifikace} - přiřazení konkrétního skutkového děje ke konkrétnímu paragrafu, odstavci a kvalifikačnímu znaku.
\end{itemize}
Jeden útok se proto popíše různě. Technicky řeknu "ransomware", incidentně "škodlivý kód + narušení dostupnosti + únik dat" a právně budu řešit neoprávněný přístup, zásah do dat, vydírání, případně podvod, porušení tajemství zpráv a neoprávněné nakládání s osobními údaji.

\subsection*{Cyber-dependent, cyber-enabled, cyber-supported}
Základní rozlišení podle role technologie:
\begin{itemize}
  \item \textbf{cyber-dependent} - bez počítače nebo sítě by trestný čin nedával smysl: neoprávněný přístup, malware, DDoS, poškození dat,
  \item \textbf{cyber-enabled} - klasická kriminalita, kterou technologie usnadňuje nebo škáluje: podvod, vydírání, stalking, krádež identity, šíření obsahu, porušení autorských práv,
  \item \textbf{cyber-supported} - technologie je podpůrný prostředek: komunikace pachatelů přes šifrované aplikace, ukládání dokumentů v cloudu, vyhledávání informací o cíli.
\end{itemize}
Toto dělení není samo o sobě právní kvalifikací, ale pomáhá ve zkouškové odpovědi. Ransomware je cyber-dependent v části průniku a zašifrování dat, ale cyber-enabled v části vydírání a vymáhání výkupného.

\subsection*{Prameny práva}
\subsubsection*{Národní prameny}
Národně je základem \textbf{trestní zákoník} a \textbf{trestní řád}. Trestní zákoník obsahuje skutkové podstaty, trestní řád dává procesní postupy pro dokazování, zajištění věcí, dat, prohlídky, odposlechy a mezinárodní součinnost.

Trojice typicky "počítačových" skutkových podstat:
\begin{itemize}
  \item \textbf{§ 230 TZ - neoprávněný přístup k počítačovému systému a nosiči informací}. Odst. 1 míří na překonání bezpečnostního opatření a získání přístupu. Odst. 2 míří na neoprávněné užití dat, poškození, změnu, vymazání, potlačení dat nebo zásah do vybavení.
  \item \textbf{§ 231 TZ - opatření a přechovávání přístupového zařízení a hesla}. Postihuje přípravné jednání: opatření, držení, zpřístupnění nebo výrobu nástroje, hesla či kódu s úmyslem spáchat trestný čin podle § 230. Není trestné samotné držení administrátorského nebo penetračního nástroje bez trestního úmyslu.
  \item \textbf{§ 232 TZ - poškození záznamu v počítačovém systému a na nosiči informací z nedbalosti}. Je důležitý jako nedbalostní podstata; typicky se řeší hrubá nedbalost a zákonem požadovaný následek.
\end{itemize}

Další časté skutkové podstaty u cyber-enabled jednání:
\begin{itemize}
  \item \textbf{§ 209 TZ - podvod}: phishing, BEC, falešné investice, romantické podvody,
  \item \textbf{§ 175 TZ - vydírání}: ransomware a hrozba zveřejněním dat,
  \item \textbf{§ 234 TZ - neoprávněné opatření, padělání a pozměnění platebního prostředku}: platební karty, internetbanking, platební údaje,
  \item \textbf{§ 182 TZ - porušení tajemství dopravovaných zpráv}: neoprávněné zachycení komunikace,
  \item \textbf{§ 192 a § 193 TZ}: materiály zachycující sexuální zneužívání dětí,
  \item \textbf{§ 270 TZ}: porušování autorských práv,
  \item \textbf{§ 355 a § 356 TZ}: nenávistné projevy a podněcování nenávisti.
\end{itemize}

\pozor{U § 230 odst. 1 nestačí říct "někdo se podíval na data". Česká úprava vyžaduje překonání bezpečnostního opatření. Pokud někdo data přímo změní, smaže nebo poškodí, řeší se zejména § 230 odst. 2, kde je logika jiná.}

\subsubsection*{Evropské prameny}
Na úrovni EU je důležitá směrnice 2013/40/EU o útocích na informační systémy. Harmonizuje nezákonný přístup, zásahy do systému a dat, nezákonné zachycování a nástroje k páchání útoků. Neříká českému soudu přímo, koho odsoudit; zavazuje státy, aby měly odpovídající skutkové podstaty.

U platebních podvodů je důležitá směrnice (EU) 2019/713 o potírání podvodů v oblasti bezhotovostních platebních prostředků. V ČR se promítla mimo jiné do změn kolem § 234 TZ. Typicky se hodí u phishingu na karty, internetbanking, platební údaje a elektronické peníze.

Procesně jsou důležité evropský vyšetřovací příkaz podle směrnice 2014/41/EU, Eurojust, Europol/EC3 a e-evidence balíček: nařízení 2023/1543 a směrnice 2023/1544. E-evidence zavede od 18. 8. 2026 EPOC a EPOC-PR pro přeshraniční získávání uložených elektronických důkazů přímo od poskytovatelů. K 16. 6. 2026 je tedy přesné říct, že jde o platný, ale teprve budoucí použitelný režim.

\subsubsection*{Mezinárodní prameny}
Základ je \textbf{Budapešťská úmluva o kyberkriminalitě}. Je to nejdůležitější funkční mezinárodní rámec: harmonizuje hmotné trestní právo, stanoví procesní nástroje pro elektronické důkazy a vytváří pravidla spolupráce. U hmotného práva přibližně odpovídá: nezákonný přístup - § 230 odst. 1, zásah do dat a systému - § 230 odst. 2, zneužití zařízení a hesel - § 231, zachycování komunikace - ochrana tajemství dopravovaných zpráv.

Druhý dodatkový protokol k Budapešťské úmluvě má posílit přímější přeshraniční získávání údajů, spolupráci s poskytovateli, registrátory domén, nouzovou spolupráci, videokonference a společná vyšetřování. K 16. 6. 2026 má podle Rady Evropy desítky signatářů, ale jen čtyři ratifikace; protože je potřeba pět ratifikací, není ještě v platnosti.

Úmluva OSN proti kyberkriminalitě také zatím není v platnosti; UNTC uvádí 77 signatářů a 3 strany, zatímco vstup vyžaduje 40 ratifikací. U zkoušky stačí říct: Budapešť je funkční rámec, 2. protokol a OSN úmluva jsou zatím neúčinné, podpis není ratifikace.

\subsection*{Typická trestná činnost}
\begin{itemize}
  \item \textbf{Hacking} - neoprávněný přístup do systému; často navazuje zneužití dat.
  \item \textbf{Malware a ransomware} - vytvoření, šíření nebo použití škodlivého kódu; poškození dat; vydírání; dvojí vydírání při exfiltraci a hrozbě zveřejněním.
  \item \textbf{Phishing} - podvodné získání údajů; často podvod, neoprávněné opatření platebního prostředku, zneužití identity.
  \item \textbf{DDoS} - omezení dostupnosti služby; podle okolností poškození systému, vydírání, útok na službu.
  \item \textbf{Krádež a prodej databáze} - neoprávněný přístup, nakládání s daty, osobní údaje, obchodní tajemství.
  \item \textbf{Online vydírání a sextortion} - vydírání, případně trestné činy proti důstojnosti a soukromí.
  \item \textbf{Nenávistné a extremistické projevy online} - podle obsahu podpora a propagace hnutí, podněcování nenávisti atd.
  \item \textbf{Nelegální tržiště a kryptoměny} - prodej přístupů, malware, drogy, legalizace výnosů, platební infrastruktura pachatelů.
  \item \textbf{Autorskoprávní kriminalita} - P2P, nelegální streamy, warez, prodej licenčních klíčů.
\end{itemize}

\subsection*{Právní kvalifikace}
Právní kvalifikace znamená převést skutkový děj do právního jazyka: určit, zda jednání naplňuje znaky konkrétního trestného činu. Nestačí říct "je to hacking" nebo "je to phishing". Musím rozebrat:
\begin{itemize}
  \item \textbf{objekt} - chráněný zájem, například integrita dat, důvěrnost komunikace, majetek,
  \item \textbf{objektivní stránku} - jednání, následek a příčinnou souvislost,
  \item \textbf{subjekt} - pachatel a jeho případné zvláštní vlastnosti,
  \item \textbf{subjektivní stránku} - zavinění, tedy úmysl nebo nedbalost.
\end{itemize}

U kyberkriminality je běžný \textbf{souběh} více trestných činů. Jeden ransomware může být § 230 odst. 1, § 230 odst. 2, § 175 vydírání, případně § 209 podvod, § 182 porušení tajemství dopravovaných zpráv a další podstaty. Důležitá jsou i stadia: příprava, pokus a dokonání. U § 231 se trestá už přípravná logika nástrojů a přístupů, ale jen s účelovým úmyslem.

Kvalifikační znaky zvyšují sazbu: škoda, rozsah, organizovaná skupina, získaný prospěch, zvlášť chráněný systém, větší počet poškozených nebo závažný dopad. Hranice škody podle § 138 TZ: větší škoda alespoň 100 tis. Kč, značná škoda alespoň 1 mil. Kč, škoda velkého rozsahu alespoň 10 mil. Kč.

\priklad{Phishing na internetbanking: vytvoření falešné stránky je příprava nebo prostředek, zaslání e-mailů je sociální inženýrství, zadání údajů obětí vede ke zneužití platebního prostředku a převod peněz je typicky podvod nebo související majetková kriminalita.}

\subsection*{Související postupy a kritéria}
Postup u kyberkriminality je prakticky důležitý, protože elektronické důkazy rychle mizí a často leží u poskytovatelů v zahraničí.

\subsubsection*{Typická procesní osa}
\begin{enumerate}
  \item \textbf{Podnět a prvotní informace} - trestní oznámení, oznámení incidentu, poznatek policie, banky, CSIRT týmu nebo zahraničního partnera.
  \item \textbf{Prověřování} - policie zjišťuje, zda došlo k trestnému činu, kdo je poškozený, jaký je vektor útoku, jaká data chybí a co je třeba okamžitě zajistit.
  \item \textbf{Elektronické místo činu} - logy, účty, cloud, servery, domény, kryptopeněženky, komunikace, IoC. Často se nejdřív řeší uchování dat, aby nezmizela.
  \item \textbf{Fyzické místo činu} - domovní prohlídka, prohlídka jiných prostor, zajištění počítačů, telefonů, disků, seed frází, hardwarových peněženek a poznámek.
  \item \textbf{Zpracování důkazů} - forenzní kopie, hash, analýza, znalecký posudek, rekonstrukce časové osy.
  \item \textbf{Právní kvalifikace a stíhání} - státní zástupce a policie posuzují skutkové podstaty, škodu, zavinění, souběh a mezinárodní prvky.
  \item \textbf{Dokazování před soudem} - soud hodnotí zákonnost, věrohodnost, integritu a význam důkazů.
\end{enumerate}

\subsubsection*{Procesní nástroje}
OČTŘ pracují hlavně s vydáním a odnětím věci, domovní prohlídkou, prohlídkou jiných prostor, uchováním dat, vyžádáním dat od držitele, provozními a lokalizačními údaji, odposlechem budoucí komunikace, znaleckým zkoumáním a mezinárodními žádostmi. Detailněji je to v otázce 21, ale u otázky 20 je nutné ukázat, že kvalifikace a proces spolu souvisejí: bez správně zajištěných důkazů nelze skutkový děj právně prokázat.

Důležitá kritéria:
\begin{itemize}
  \item \textbf{zákonnost} - existuje právní titul a správný procesní postup,
  \item \textbf{nezbytnost a proporcionalita} - zásah do soukromí odpovídá závažnosti věci,
  \item \textbf{rychlost} - logy a data se mažou, účty se ruší, kryptoměny se přesouvají,
  \item \textbf{integrita} - důkaz se nesmí nekontrolovaně měnit,
  \item \textbf{chain of custody} - kdo, kdy a proč s důkazem nakládal,
  \item \textbf{jurisdikce} - kde jsou data, poskytovatel, pachatel, oběť a dopad,
  \item \textbf{použitelnost důkazu} - i technicky dobrý důkaz může být procesně problematický, pokud byl získán nezákonně.
\end{itemize}

\subsection*{Mezinárodní spolupráce}
Kyberkriminalita je typicky přeshraniční: oběť je v ČR, server v Irsku, pachatel jinde, doména registrovaná přes další stát, platba v kryptoměnách a komunikace přes globální platformu. Problémem je jurisdikce, rychlost, rozdílná trestnost, lokalizace dat, ochrana práv a použitelnost důkazů.

Používají se:
\begin{itemize}
  \item \textbf{vzájemná právní pomoc (MLA)} - formální žádost jednoho státu druhému, pomalejší, ale klasická,
  \item \textbf{evropský vyšetřovací příkaz (EVP/EIO)} - v EU orgán žádá orgán jiného členského státu o vyšetřovací úkon; funguje na principu vzájemného uznávání,
  \item \textbf{Europol/EC3, J-CAT a Eurojust} - operativní koordinace, analytická podpora, společné akce, koordinace státních zástupců,
  \item \textbf{síť 24/7 podle Budapešťské úmluvy} - rychlý kontakt pro uchování a směrování žádostí,
  \item \textbf{EPOC a EPOC-PR} - od 18. 8. 2026 přímější režim vůči poskytovatelům v EU,
  \item \textbf{dobrovolná spolupráce poskytovatelů} - možná v mezích práva, pravidel poskytovatele a typu dat; nenahrazuje procesní titul tam, kde je povinný.
\end{itemize}

\pravobox{E-evidence v otázce 20: k 16. 6. 2026 jde o budoucí režim použitelný od 18. 8. 2026. EPOC je evropský příkaz k vydání uložených elektronických důkazů; standardní lhůta je 10 dnů, v naléhavém případě 8 hodin. EPOC-PR je evropský příkaz k uchování dat; data se uchovají 60 dnů s možností prodloužení o 30 dnů. Režim pracuje se čtyřmi kategoriemi: údaje o účastníkovi/službě, údaje sloužící výlučně k identifikaci uživatele, provozní údaje a obsahová data. Příkazy vydává nebo validuje soud, státní zástupce či jiný příslušný orgán podle typu dat; u provozních a obsahových dat jsou požadavky přísnější. Existují důvody odmítnutí a přezkumu, například základní práva, imunity a výsady, profesní tajemství nebo zjevné porušení Listiny EU. EPOC není živý odposlech, neprolamuje end-to-end šifrování a nevytvoří data, která poskytovatel nemá. Dánsko režimem není vázáno.}

\subsection*{Modelové situace}
\begin{enumerate}
  \item \textbf{Hacking a krádež databáze}: § 230 odst. 1 za přístup po překonání zabezpečení, § 230 odst. 2 za užití nebo exfiltraci dat; podle rozsahu škoda, osobní údaje, obchodní tajemství, GDPR breach a případně ZKB incident.
  \item \textbf{Ransomware s výkupným}: průnik, zašifrování dat, zásah do dostupnosti, vydírání; při exfiltraci dvojí vydírání; typicky organizovaná a přeshraniční skupina.
  \item \textbf{Phishing na internetbanking}: podvod, platební prostředek, případně neoprávněný přístup, pokud se pachatel přihlásí do účtu oběti.
  \item \textbf{Prodej přístupových údajů na fóru}: § 231, pokud je dán účelový úmysl spáchat § 230; není nutné, aby už došlo k průniku.
  \item \textbf{Admin z hrubé nedbalosti smaže produkční databázi}: může jít o § 232, ale jen při naplnění hrubé nedbalosti a zákonného následku; jinak spíš civilní, pracovní nebo správní rovina.
  \item \textbf{Důkazy u poskytovatele v Irsku}: podle času a režimu MLA/EVP, od 18. 8. 2026 e-evidence EPOC/EPOC-PR přímo vůči poskytovateli nebo jeho určenému adresátovi.
\end{enumerate}

\zkouska{U kvalifikace vždy postupuj od faktů: co pachatel udělal, k jakému systému nebo datům, s jakým úmyslem, jaká škoda a jaký následek. Nepiš "je to hacking", ale převeď hacking na skutkové podstaty a řekni, jakými procesními kroky to lze prokázat.}

\zdroje{Trestní zákoník; trestní řád; směrnice 2013/40/EU; směrnice 2019/713; směrnice 2014/41/EU o EVP; Budapešťská úmluva; e-evidence nařízení 2023/1543 a směrnice 2023/1544; UNTC status Úmluvy OSN proti kyberkriminalitě; referenční text z přiloženého PDF a starší výpisky.}

\newpage

\otazka{21. Elektronické důkazy a jejich zajišťování}{Procesní instituty a jejich praktické využití, nakládání s elektronickými důkazy, elektronické dokumenty.}{BVV03K}

\subsection*{Kostra odpovědi}
\begin{enumerate}
  \item Definovat elektronický důkaz.
  \item Vysvětlit jeho vlastnosti: volatilita, objem, přeshraničnost, metadata, šifrování.
  \item Popsat procesní instituty v trestním řádu: § 78, § 79, § 82 a násl., § 7b, § 88a, § 88.
  \item Doplnit zákon o elektronických komunikacích: metadata vs obsah, role operátorů a ČTÚ.
  \item Vysvětlit přeshraniční režimy: MLA, EVP, Budapešť, EPOC/EPOC-PR.
  \item Vysvětlit forenzní zásady: live/dead forenzika, integrita, dokumentace, hash, chain of custody.
  \item Popsat elektronické dokumenty, podpisy, pečetě a časová razítka.
\end{enumerate}

\subsection*{Co je elektronický důkaz}
Elektronický důkaz je jakákoli digitální informace použitelná k prokázání skutečnosti v řízení. Může jít o soubor, e-mail, log, metadata, databázi, fotografii, zprávu v aplikaci, záznam přihlášení, IP adresu, cookies, data z telefonu, cloudový účet, kamerový záznam, GPS stopu, síťový provoz, kryptoměnovou transakci nebo konfiguraci systému.

Elektronické důkazy jsou zvláštní tím, že jsou snadno změnitelné, volatilní, často u třetích stran, šifrované, objemné a přeshraniční. Obsah RAM zmizí vypnutím, logy rotují, cloudová data mohou být v jiné zemi a časové údaje mohou být v různých časových pásmech. Proto nestačí důkaz "mít"; je potřeba prokázat, že byl získán zákonně, že nebyl změněn a že interpretace dává smysl.

Zkouškově je klíčové rozlišit \textbf{fyzický nosič} a \textbf{data}. Zajištění notebooku nebo telefonu neznamená automaticky neomezené oprávnění prohledat celý digitální život člověka. Jiný zásah je zajistit zařízení, jiný vytvořit forenzní kopii, jiný prohledávat obsah, jiný vyžádat data u poskytovatele a jiný zachycovat budoucí komunikaci. Stejně tak je rozdíl mezi \textbf{historicky uloženými daty} a \textbf{živou nebo budoucí komunikací}: starý e-mail uložený ve schránce se nezískává jako odposlech, zatímco budoucí obsah komunikace už spadá do přísnějšího režimu odposlechu.

Při zajišťování se vždy řeší rozsah a proporcionalita. Příkaz má být co nejpřesnější podle účelu řízení, aby se z trestního zajištění nestalo plošné prohledání všeho. Zvláštní opatrnost vyžadují privilegovaná data, zejména komunikace s obhájcem nebo advokátní tajemství, zdravotnická dokumentace, novinářské zdroje a soukromí třetích osob. V praxi proto pomáhá filtrovat data, dokumentovat kroky, oddělit irelevantní obsah a vždy umět vysvětlit, proč byl zásah nezbytný.

\pozor{Elektronický důkaz není automaticky slabší než listina. Slabý je tehdy, když neumím prokázat původ, integritu, čas, zákonnost získání nebo souvislost s dokazovanou skutečností.}

\subsection*{Procesní instituty trestního řádu}
Je potřeba rozlišovat, co chci získat: zařízení, uložená data, metadata, historický obsah u poskytovatele, nebo budoucí komunikaci.
\begin{itemize}
  \item \textbf{§ 78 TŘ - vydání věci} - osoba vydá věc důležitou pro trestní řízení; může jít o telefon, disk, notebook, token, dokumentaci nebo nosič.
  \item \textbf{§ 79 TŘ - odnětí věci} - když věc není vydána dobrovolně, lze ji odňat procesním postupem.
  \item \textbf{§ 82 a násl. TŘ - domovní prohlídka a prohlídka jiných prostor} - typicky zajištění techniky, serverů, listinných poznámek, seed frází, hardwarových peněženek a přístupových údajů.
  \item \textbf{§ 7b TŘ - uchování dat} - rychlé "zmrazení" dat, aby nezmizela dřív, než se opatří další procesní titul k vydání; věcně podobné quick freeze logice Budapešťské úmluvy.
  \item \textbf{§ 88a TŘ - provozní a lokalizační údaje} - metadata: kdo s kým, kdy, jak dlouho, odkud; ne obsah komunikace.
  \item \textbf{§ 88 TŘ - odposlech a záznam telekomunikačního provozu} - obsah budoucí komunikace, jen za přísných podmínek a zpravidla na příkaz soudu.
  \item \textbf{znalecké zkoumání} - forenzní analýza, interpretace stop, posouzení malwaru, časové osy, zařízení nebo logů.
\end{itemize}

\zkouska{§ 88 = obsah budoucí komunikace. § 88a = metadata. § 7b = uchování existujících dat. § 78/79 = vydání nebo odnětí věci/dat. Tohle rozlišení je u otázky 21 klíčové.}

\subsection*{Zákon o elektronických komunikacích}
Zákon o elektronických komunikacích (č. 127/2005 Sb.) je důležitý proto, že velká část stop leží u operátorů. Upravuje sítě a služby elektronických komunikací, chrání důvěrnost komunikace, stanoví povinnosti součinnosti a dohledovou roli Českého telekomunikačního úřadu.

Nejdůležitější je rozdíl mezi \textbf{metadaty} a \textbf{obsahem}. Metadata jsou provozní a lokalizační údaje: kdo komunikoval, kdy, odkud, jak dlouho a přes jakou službu. Obsah je text zprávy, zvuk hovoru nebo samotný obsah komunikace. Operátor běžně uchovává metadata v zákonném režimu, ale neukládá obsah minulých hovorů jen proto, aby ho později vydal policii. Naopak cloudová nebo e-mailová služba může mít historický obsah uložený na serveru; ten se pak získává vydáním/uchováním dat, ne odposlechem.

\pozor{Zákon o elektronických komunikacích sám o sobě policii "nedává odposlech". Stanoví povinnosti operátorů a součinnost, ale procesní oprávnění musí přijít z trestního řádu nebo jiného zvláštního zákona.}

\subsection*{Přeshraniční elektronické důkazy}
U cloudu, e-mailu, sociálních sítí a globálních poskytovatelů je otázka často přeshraniční. Používají se:
\begin{itemize}
  \item \textbf{vzájemná právní pomoc (MLA)} - klasická žádost mezi státy,
  \item \textbf{evropský vyšetřovací příkaz (EVP/EIO)} - v EU orgán žádá orgán druhého státu o zajištění nebo předání důkazu,
  \item \textbf{Budapešťská úmluva} - zejména rychlé uchování uložených dat a provozních údajů a síť 24/7,
  \item \textbf{Europol/Eurojust} - koordinace, analytická a justiční podpora,
  \item \textbf{e-evidence} - EPOC a EPOC-PR od 18. 8. 2026.
\end{itemize}

\pravobox{EPOC/EPOC-PR v otázce 21: nařízení 2023/1543 zavádí evropský vydávací příkaz a evropský uchovávací příkaz pro elektronické důkazy v trestních věcech; k 16. 6. 2026 jde stále o budoucí režim použitelný od 18. 8. 2026. EPOC míří na vydání uložených dat, standardně do 10 dnů a v naléhavém případě do 8 hodin. EPOC-PR míří na uchování dat bez odkladu po dobu 60 dnů, s možností prodloužení o 30 dnů. Směrnice 2023/1544 řeší určené provozovny a právní zástupce poskytovatelů, aby bylo komu příkazy doručovat. Data se dělí na údaje o účastníkovi/službě, údaje sloužící výlučně k identifikaci uživatele, provozní údaje a obsah. Čím citlivější kategorie, tím přísnější vydání nebo validace příkazu. Režim chrání základní práva, profesní tajemství a umožňuje odmítnutí nebo přezkum v zákonných situacích. Nevztahuje se na Dánsko, neznamená živý odposlech a neprolamuje šifrování.}

\pozor{EVP a EPOC nejsou totéž. EVP je model "orgán žádá orgán" a může se týkat různých vyšetřovacích úkonů. EPOC je model "orgán míří přímo na poskytovatele nebo jeho určeného adresáta" a týká se uložených elektronických důkazů.}

\subsection*{Nakládání s elektronickými důkazy}
Forenzní postup má dva cíle: zajistit relevantní data a zároveň udržet jejich důkazní hodnotu. Základní zásady:
\begin{itemize}
  \item \textbf{zdokumentovat výchozí stav} - fotografie, zapojení, čas, obrazovka, síťové připojení, kdo byl přítomen,
  \item \textbf{rozhodnout live vs dead forenziku} - běžící šifrovaný systém může vyžadovat RAM a síťová spojení; jindy je lepší řízené vypnutí,
  \item \textbf{minimalizovat změny} - pokud možno analyzovat kopii, originál zapečetit,
  \item \textbf{vytvořit forenzní obraz} - bitová kopie zařízení nebo relevantních dat; u klasických nosičů ideálně přes write-blocker,
  \item \textbf{hashovat} - například SHA-256 originálu a kopie, aby bylo možné prokázat integritu kopie,
  \item \textbf{dokumentovat nástroje a metodu} - název, verze, nastavení, časové pásmo, postup,
  \item \textbf{chain of custody} - nepřerušený řetězec držení: kdo, kdy, kde, proč a jak s důkazem nakládal,
  \item \textbf{respektovat rozsah příkazu} - nepřiměřené "prohledání všeho" může narazit na zákonnost, soukromí, advokátní tajemství nebo privilegovaná data.
\end{itemize}

Hash prokazuje, že kopie odpovídá tomu, co bylo zajištěno. Nedokazuje ale pravdivost obsahu, autenticitu původce ani to, že data nebyla zmanipulována před zajištěním. To se prokazuje dalšími důkazy, kontextem, logy, podpisy, svědeckými výpověďmi a znaleckým posouzením.

\priklad{Běžící server je šifrovaný. Pokud ho vyšetřovatel bez rozmyslu vypne, může přijít o klíče v paměti a data už neotevře. Někdy je proto vhodná živá akvizice: zajistit RAM, síťová spojení a běžící procesy, potom pořídit obraz disku, spočítat hash a analyzovat kopii.}

\subsection*{Co důkaz znehodnotí}
\begin{itemize}
  \item rychlé otevření originálu "jen pro kontrolu", protože se změní metadata,
  \item chybějící hash nebo nejasný postup kopírování,
  \item díra v chain of custody,
  \item nezdokumentovaný zásah specialisty,
  \item nesynchronizovaný čas nebo chybějící časové pásmo,
  \item nevalidovaná metoda bez možnosti opakování,
  \item překročení rozsahu příkazu nebo nepřiměřený zásah do soukromí,
  \item směšování incident response a trestního zajištění důkazu bez jasného záznamu.
\end{itemize}

\subsection*{Modelové situace}
\begin{enumerate}
  \item \textbf{Ransomware na běžícím šifrovaném serveru}: řešit live forenziku, RAM, síťová spojení, IoC, logy ze SIEM, potom forenzní obraz, hash a analýza kopie.
  \item \textbf{Mobil podezřelého}: izolace od sítě kvůli riziku vzdáleného smazání, ale postup volí specialista; Faradayův obal nebo režim letadlo nejsou magická univerzální řešení.
  \item \textbf{Data v cloudu}: nejprve uchování dat, potom vydání dat správným procesním titulem; přeshraničně MLA/EVP/EPOC podle situace.
  \item \textbf{Metadata od operátora}: § 88a, ne § 88; jde o provozní a lokalizační údaje, ne obsah komunikace.
  \item \textbf{Historický e-mail}: není to odposlech budoucí komunikace; jde o uložený obsah u poskytovatele nebo v zařízení.
\end{enumerate}

\subsection*{Civilní a správní řízení}
Elektronické důkazy se nepoužívají jen v trestním řízení. V civilním sporu může být důkazem e-mail, elektronická smlouva, log z informačního systému, datová zpráva nebo elektronicky podepsaný dokument. Ve správním řízení může jít o spisovou službu, datové schránky, kamerové záznamy nebo dokumenty ve veřejné správě. Důležitý rozdíl je, že civilní nebo správní orgán nemá automaticky trestněprocesní donucovací nástroje typu odposlechu, domovní prohlídky nebo zajištění dat u poskytovatele podle trestního řádu. I tam se ale hodnotí autenticita, integrita, původ, čas a zákonnost získání.

\subsection*{Elektronické dokumenty}
Elektronický dokument není méněcenný jen proto, že je elektronický. eIDAS stanoví, že elektronickému dokumentu a elektronickému podpisu nelze upírat právní účinky pouze kvůli elektronické formě. Důkazní hodnota ale závisí na tom, zda lze prokázat autenticitu, integritu a čas.

Elektronický podpis dokládá vazbu osoby k dokumentu a integritu od podpisu. Elektronická pečeť dokládá původ u právnické osoby nebo instituce. Časové razítko dokládá, že data existovala v určité podobě v určitém čase. Ani podpis, ani pečeť, ani razítko nedokazují pravdivost obsahu. Dokazují pouze původ, integritu a časové souvislosti.

U dokumentu jako důkazu řeším tři otázky: \textbf{autenticita} - od koho dokument pochází, \textbf{integrita} - zda se nezměnil, \textbf{čas} - kdy existoval v dané podobě. Kvalifikovaný podpis má účinky vlastnoručního podpisu, kvalifikovaná pečeť zakládá domněnku integrity a původu od právnické osoby a kvalifikované časové razítko podporuje důkaz o existenci dat v čase. U dlouhodobé archivace je potřeba myslet na stárnutí certifikátů a algoritmů, proto se používají validační data a navazující časová razítka.

\pozor{Kvalifikovaný podpis má účinky vlastnoručního podpisu, ale neznamená to, že obsah dokumentu je pravdivý. Podepsat lze i nepravdivé prohlášení.}

\zkouska{U elektronických důkazů propoj právo a techniku: procesní titul podle trestního řádu, rozdíl metadata/obsah, přeshraniční nástroj, a potom forenzní integrita přes kopii, hash, dokumentaci a chain of custody.}

\zdroje{Trestní řád (§ 7b, § 78, § 79, § 82 a násl., § 88, § 88a); zákon o elektronických komunikacích (č. 127/2005 Sb.); směrnice 2014/41/EU; Budapešťská úmluva čl. 16, 17, 29, 30, 35; e-evidence nařízení 2023/1543 a směrnice 2023/1544; eIDAS; ISO/IEC 27037; přiložený referenční text a starší výpisky.}

\newpage

\otazka{22. Ochrana osobních údajů}{Právní úprava. Principy a zásady zpracování osobních údajů - definice, základní zásady, posouzení rizik, test proporcionality. Účely zpracování. Právní tituly. Zákon o zpracování osobních údajů - ÚOOÚ a jeho úloha a postavení v České republice.}{BI301K}

\subsection*{Kostra odpovědi}
\begin{enumerate}
  \item Právní rámec: GDPR jako přímo použitelné nařízení, zákon o zpracování osobních údajů, zvláštní předpisy a ePrivacy/obchodní sdělení.
  \item Pojmy: osobní údaj a identifikovatelnost, subjekt údajů, správce, zpracovatel, příjemce, pseudonymizace vs anonymizace.
  \item Zvláštní kategorie údajů: čl. 6 a čl. 9 zároveň, čl. 10 pro trestní údaje.
  \item Zásady čl. 5 na příkladech: zákonnost, transparentnost, účel, minimalizace, přesnost, omezení uložení, integrita/důvěrnost, odpovědnost.
  \item Nejdřív účel, potom právní titul: šest titulů čl. 6, souhlas, soft opt-in a oprávněný zájem.
  \item Oprávněný zájem: LIA/balancing test - oprávněný cíl, nezbytnost, vyvážení zájmů a právo námitky.
  \item DPIA, proporcionalita, privacy by design/default a DPO.
  \item Práva subjektů: přístup, oprava, výmaz, omezení, přenositelnost, námitka, automatizované rozhodování.
  \item Porušení zabezpečení, ÚOOÚ, EDPB, one-stop-shop a předávání mimo EHP.
\end{enumerate}

\subsection*{Právní úprava}
Základem je GDPR, tedy nařízení EU 2016/679, použitelné od 25. května 2018. Je přímo použitelné a sjednocuje pravidla ochrany osobních údajů v EU. To je rozdíl proti směrnicím typu NIS2: GDPR se netransponuje, ale přímo platí. Český zákon o zpracování osobních údajů (č. 110/2019 Sb.) je adaptační zákon. GDPR nenahrazuje, jen doplňuje tam, kde nařízení nechává prostor členskému státu.

Český zákon řeší zejména postavení ÚOOÚ, některé výjimky pro žurnalistiku, vědu, archivnictví a svobodu projevu, samostatný režim pro zpracování v trestněprávní oblasti a také věkovou hranici dítěte pro souhlas u služeb informační společnosti: v ČR 15 let. Vedle toho existují zvláštní předpisy pro zdravotnictví, pracovní právo, elektronické komunikace, obchodní sdělení, bankovnictví, školství, archivnictví, policii a zpravodajské služby.

GDPR není "zákaz zpracování dat". Je to pravidlový rámec: osobní údaje zpracovávat lze, ale musí být jasný účel, právní titul, přiměřenost, bezpečnost a odpovědnost.

\subsubsection*{GDPR, trestněprávní oblast a územní působnost}
GDPR se neuplatní na všechno zpracování osobních údajů veřejnou mocí. Pokud příslušné orgány zpracovávají osobní údaje pro účely prevence, vyšetřování, odhalování nebo stíhání trestných činů, výkonu trestů nebo ochrany před hrozbami pro veřejnou bezpečnost, nastupuje zvláštní režim směrnice 2016/680, tzv. Law Enforcement Directive, a český zákon o zpracování osobních údajů. Stejný orgán proto může být jednou v režimu GDPR a jindy v režimu LED podle účelu. Policie při personalistice nebo nákupu služeb řeší GDPR; při vyšetřování trestného činu řeší zvláštní trestněprávní režim.

Územní působnost GDPR není omezena jen na firmy fyzicky v EU. GDPR se použije na zpracování v souvislosti s provozovnou správce nebo zpracovatele v EU. Vedle toho dopadá i na subjekty mimo EU, pokud nabízejí zboží nebo služby osobám v EU nebo monitorují jejich chování v EU, například cíleným trackingem, profilováním nebo behaviorální reklamou. Neznamená to, že každý web na světě automaticky spadá pod GDPR; rozhoduje reálné cílení na osoby v EU nebo monitorování jejich chování.

\pozor{Domácí výjimka se týká čistě osobní nebo domácí činnosti, například soukromého adresáře nebo rodinné fotografie. Neplatí automaticky, pokud osoba zveřejňuje údaje neurčité veřejnosti na internetu, provozuje veřejnou stránku, systematicky monitoruje okolí nebo zasahuje do práv cizích osob.}

\subsection*{Základní pojmy}
\textbf{Osobní údaj} je informace o identifikované nebo identifikovatelné fyzické osobě. Identifikovatelnost znamená, že osobu lze určit přímo nebo nepřímo pomocí rozumně dostupných prostředků. Jméno a rodné číslo jsou zřejmé příklady, ale osobním údajem může být i IP adresa, cookie identifikátor, poloha, hlas, fotografie, pracovní identifikátor nebo kombinace údajů.

Ne každý technický identifikátor je osobním údajem vždy a pro každého. Dynamická IP adresa nebo cookie je osobní údaj tehdy, pokud existují rozumně pravděpodobné prostředky, jak ji spojit s osobou. Pro poskytovatele služby to často osobní údaj bude; pro subjekt bez možnosti vazby to tak být nemusí.

\textbf{Subjekt údajů} je fyzická osoba, které se údaje týkají. \textbf{Správce} určuje účely a prostředky zpracování, tedy proč a jak se údaje zpracovávají. \textbf{Zpracovatel} zpracovává údaje pro správce a podle jeho pokynů, typicky cloud, externí účetní, mailingový nástroj nebo IT dodavatel; vztah se upravuje zpracovatelskou smlouvou podle čl. 28 GDPR. \textbf{Příjemce} je ten, komu jsou údaje zpřístupněny. \textbf{Zpracování} je skoro jakákoli operace s údaji: sběr, ukládání, změna, použití, předání, nahlížení, zveřejnění nebo výmaz.

\textbf{Pseudonymizace} nahrazuje identifikátory kódem, ale klíč nebo vazba existuje odděleně, takže GDPR stále platí. \textbf{Anonymizace} musí být prakticky nevratná; osobu už nelze rozumně pravděpodobnými prostředky znovu identifikovat, vyčlenit ani propojit. Pouhé odstranění jména není anonymizace, pokud zbylá data pořád umožní re-identifikaci.

\pozor{Pseudo není anonym. Pseudonymizovaná data zůstávají osobními údaji, protože existuje vazba zpět na člověka.}

\subsection*{Zvláštní kategorie údajů}
Zvláštní kategorie podle čl. 9 GDPR jsou citlivější osobní údaje: zdravotní údaje, genetické údaje, biometrické údaje zpracovávané za účelem jedinečné identifikace, údaje o sexuální orientaci, politických názorech, náboženském nebo filozofickém přesvědčení, členství v odborech a rasovém nebo etnickém původu. Jejich zpracování je v zásadě zakázané a teprve výjimka podle čl. 9 zákaz prolomí.

U zvláštních kategorií je potřeba splnit dvě vrstvy najednou:
\begin{itemize}
  \item \textbf{čl. 6} - obecný právní titul, tedy proč osobní údaj vůbec smím zpracovat,
  \item \textbf{čl. 9} - zvláštní výjimka, která dovolí zpracovat citlivý údaj.
\end{itemize}
Splnit jen jednu vrstvu nestačí. Nemocnice například zpracovává zdravotní údaje na základě právní povinnosti nebo poskytování zdravotní péče podle čl. 6 a současně používá výjimku podle čl. 9 pro zdravotní péči. Údaje o odsouzeních a trestných činech mají vlastní režim v čl. 10.

\priklad{Fitness aplikace sbírá počet kroků a tepovou frekvenci. Pokud z toho vyvozuje zdravotní stav nebo používá data pro zdravotní doporučení, může se dostat ke zvláštní kategorii údajů. Pak nestačí "uživatel souhlasil s podmínkami"; je potřeba přesně vyřešit titul podle čl. 6 a výjimku podle čl. 9, často výslovný souhlas.}

\subsection*{Zásady zpracování}
GDPR stojí na zásadách:
\begin{itemize}
  \item \textbf{zákonnost, korektnost a transparentnost} - zpracování musí mít právní titul a člověk má srozumitelně vědět, co se děje; e-shop musí jasně informovat, že e-mail použije pro objednávku,
  \item \textbf{účelové omezení} - data se sbírají pro konkrétní účel a nelze je libovolně použít jinak; e-mail z objednávky nejde bez dalšího použít k profilování reklam,
  \item \textbf{minimalizace} - sbírat jen to, co je nezbytné; pro newsletter stačí e-mail, ne rodné číslo,
  \item \textbf{přesnost} - údaje mají být aktuální a správné; špatnou doručovací adresu je třeba opravit,
  \item \textbf{omezení uložení} - neuchovávat déle, než je potřeba; faktury podle účetních lhůt, marketingový kontakt jen do odvolání souhlasu nebo námitky,
  \item \textbf{integrita a důvěrnost} - zajistit bezpečnost: řízení přístupů, šifrování, zálohy, logování,
  \item \textbf{odpovědnost} - správce musí soulad nejen dodržet, ale i doložit: záznamy o činnostech, testy, DPIA, smlouvy, interní pravidla.
\end{itemize}

\subsection*{Účely a právní tituly}
Účel říká, proč údaje zpracovávám. Právní titul říká, proč to smím. Správné pořadí je: nejdřív účel, potom právní titul, potom rozsah údajů, doba uložení, informace subjektu a zabezpečení. Titul se nemá vybírat dodatečně jako náplast, když se původní důvod nehodí.

Čl. 6 GDPR zná šest titulů:
\begin{enumerate}
  \item \textbf{souhlas} - newsletter pro nového zájemce, pokud není jiný titul,
  \item \textbf{plnění smlouvy} - jméno a adresa pro doručení objednávky,
  \item \textbf{právní povinnost} - uchování účetních a daňových dokladů,
  \item \textbf{ochrana životně důležitých zájmů} - péče o pacienta v bezvědomí,
  \item \textbf{úkol ve veřejném zájmu nebo výkon veřejné moci} - matrika, veřejná správa,
  \item \textbf{oprávněný zájem} - bezpečnostní monitoring, ochrana majetku, vymáhání pohledávky, některý přímý marketing.
\end{enumerate}

Souhlas není univerzální náplast. Musí být svobodný, konkrétní, informovaný, jednoznačný a odvolatelný stejně snadno, jako byl udělen. Problém vzniká, když je služba podmíněna souhlasem se zpracováním, které pro službu není nezbytné: například nákup v e-shopu podmíněný souhlasem s marketingem.

U přímého elektronického marketingu je vedle souhlasu důležitý i český režim obchodních sdělení podle zákona o některých službách informační společnosti (č. 480/2004 Sb.) a tzv. soft opt-in: stávajícímu zákazníkovi lze za splnění podmínek nabízet vlastní obdobné výrobky nebo služby, pokud měl možnost jednoduše a zdarma odmítnout při získání kontaktu i v každé další zprávě.

\subsection*{Oprávněný zájem a balancing test}
Oprávněný zájem podle čl. 6 odst. 1 písm. f GDPR není volná karta. Správce musí před zpracováním projít a zdokumentovat test oprávněného zájmu, často označovaný jako LIA. Má tři kroky:
\begin{enumerate}
  \item \textbf{Test účelu / oprávněnosti cíle} - je zájem konkrétní, skutečný a zákonný? Typicky ochrana majetku, bezpečnost sítě, prevence podvodů, vymáhání pohledávky, přímý marketing vůči vlastním zákazníkům.
  \item \textbf{Test nezbytnosti} - je zpracování pro daný cíl skutečně potřebné a nejde stejného výsledku dosáhnout méně invazivně? Kamera na vstup může být nezbytná, kamera na celou jídelnu nebo šatnu už ne.
  \item \textbf{Balancing test / vyvážení} - nepřevažují práva, svobody a rozumná očekávání subjektu nad zájmem správce? Hodnotí se rozsah údajů, citlivost, vztah se subjektem, možnost námitek, dopad, transparentnost, záruky a bezpečnostní opatření.
\end{enumerate}

Výsledek má být písemně doložen, protože správce nese odpovědnost. Test oprávněného zájmu není totéž co DPIA a není totéž co ústavní test proporcionality. LIA řeší, zda mohu použít konkrétní právní titul. DPIA řeší rizika rizikového zpracování jako celek. Ústavní test proporcionality je obecnější test zásahu do základních práv, typicky u veřejné moci.

\pravobox{TEST OPRÁVNĚNÉHO ZÁJMU A JEHO HRANICE: Firma chce nasadit kameru u vstupu do skladu kvůli krádežím. Cíl je oprávněný: ochrana majetku a bezpečnost. Nezbytnost: kamera míří jen na vstup, nikoli na odpočinkovou místnost, záznam se drží krátce a přístup má jen omezený počet osob. Vyvážení: zaměstnanci a návštěvníci jsou jasně informováni, nejde o skryté sledování, záznam se používá jen při incidentu. Oprávněný zájem pravděpodobně obstojí. Kdyby kamera snímala celý pracovní prostor se zvukem a záznam se používal k výkonovému hodnocení zaměstnanců, vyvážení by nejspíš neobstálo a mohlo by být potřeba DPIA.}

\priklad{Bezpečnostní logy e-shopu: oprávněný zájem na ochraně účtů a detekci útoků je reálný. Nezbytnost znamená logovat údaje, které opravdu pomáhají detekci, ne sbírat všechno navždy. Vyvážení znamená krátké retenční lhůty, omezený přístup, informování uživatelů a nepoužívat bezpečnostní logy potichu k marketingovému profilování.}

\subsection*{Riziko, DPIA a proporcionalita}
GDPR je rizikově orientované. Čím vyšší riziko pro práva a svobody osob, tím silnější musí být opatření. Když zpracování pravděpodobně povede k vysokému riziku, provádí se \textbf{DPIA} - posouzení vlivu na ochranu osobních údajů podle čl. 35 GDPR. DPIA se dělá \textbf{před zahájením zpracování}, ne až po incidentu nebo po spuštění systému.

Typicky je DPIA potřeba u:
\begin{itemize}
  \item systematického a rozsáhlého profilování s právními nebo obdobně závažnými účinky,
  \item rozsáhlého zpracování zvláštních kategorií údajů,
  \item rozsáhlého systematického monitorování veřejně přístupných prostor,
  \item biometrické identifikace, rozsáhlého kamerového sledování, nových technologií, AI systémů nebo kombinování datových sad,
  \item zpracování, které může vést k diskriminaci, ztrátě kontroly nad daty, finanční újmě, poškození pověsti nebo omezení práv.
\end{itemize}

Obsah DPIA v praxi:
\begin{enumerate}
  \item popsat účel a zamýšlené zpracování,
  \item určit kategorie údajů, subjektů, příjemců, zdroje a dobu uložení,
  \item doložit právní titul a případně výjimku podle čl. 9,
  \item posoudit nezbytnost a přiměřenost,
  \item identifikovat rizika pro osoby,
  \item navrhnout opatření: minimalizace, pseudonymizace, šifrování, řízení přístupu, logování, omezení retence, lidský dohled,
  \item vyhodnotit zbytkové riziko.
\end{enumerate}
Pokud i po opatřeních zůstává vysoké riziko, musí správce před zahájením zpracování konzultovat ÚOOÚ.

Test proporcionality se dá vyložit jednoduše:
\begin{enumerate}
  \item \textbf{legitimní cíl} - sleduji oprávněný a jasný účel?
  \item \textbf{vhodnost} - pomůže zpracování cíle dosáhnout?
  \item \textbf{nezbytnost} - nejde to méně invazivně?
  \item \textbf{přiměřenost v užším smyslu} - nepřevažuje zásah do práv nad přínosem?
\end{enumerate}

\priklad{Kamerový systém ve škole: účel může být bezpečnost. Vhodnost ano, kamery mohou pomoci. Nezbytnost se ptá, zda nestačí méně kamer, kratší uchování nebo jen vstupy. Přiměřenost řeší, zda se nesledují šatny, toalety nebo žáci nepřiměřeně dlouho.}

\subsection*{Pověřenec, privacy by design a profilování}
\textbf{Pověřenec pro ochranu osobních údajů} (DPO) je osoba, která radí správci nebo zpracovateli, sleduje soulad s GDPR, školí, kontroluje, spolupracuje s ÚOOÚ a je kontaktním místem pro subjekty údajů. Není to "ten, kdo za všechno odpovídá" místo správce; odpovědnost za soulad má pořád správce nebo zpracovatel.

DPO je povinný zejména:
\begin{itemize}
  \item u orgánů veřejné moci a veřejných subjektů,
  \item když hlavní činnost spočívá v rozsáhlém pravidelném a systematickém monitorování osob,
  \item když hlavní činnost spočívá v rozsáhlém zpracování zvláštních kategorií údajů nebo údajů podle čl. 10.
\end{itemize}
DPO musí být odborně způsobilý, nezávislý ve výkonu své funkce, nesmí být za plnění úkolů sankcionován a nemá se dostat do střetu zájmů. Malý e-shop DPO obvykle nepotřebuje, nemocnice nebo obec ano.

\textbf{Privacy by design a by default} znamená, že ochrana údajů se zabuduje už do návrhu systému a výchozí nastavení je co nejšetrnější: sbírat méně dat, kratší retence, omezené role, auditní logy, pseudonymizace, oddělení prostředí, výchozí neveřejnost profilu.

\textbf{Profilování} je automatizované zpracování osobních údajů k hodnocení osobních aspektů člověka, například ekonomické situace, zdraví, preferencí, chování, spolehlivosti nebo pohybu. Profilování není vždy zakázané, ale u rozsáhlého a systematického profilování s významnými dopady typicky roste riziko, může být potřeba DPIA a u čistě automatizovaného rozhodování s právními nebo obdobně závažnými účinky se uplatní čl. 22 GDPR.

\priklad{Nemocnice chce nasadit AI systém pro podporu diagnostiky. Zpracovává zdravotní údaje, tedy zvláštní kategorii. Potřebuje titul podle čl. 6 a výjimku podle čl. 9, DPO je prakticky povinný kvůli rozsáhlému zpracování citlivých údajů, před spuštěním se dělá DPIA a opatření jako pseudonymizace tréninkových dat, omezení přístupů, logování a lidský dohled nad rozhodnutím.}

\subsection*{Porušení zabezpečení a práva subjektů}
Porušení zabezpečení osobních údajů je únik, ztráta, zničení, změna nebo neoprávněný přístup. Správce oznamuje dozorovému úřadu zpravidla do 72 hodin od okamžiku, kdy se o porušení dozvěděl, pokud je pravděpodobné riziko pro práva a svobody. Subjektům údajů oznamuje bez zbytečného odkladu až tehdy, pokud je vysoké riziko. Pokud je riziko nepravděpodobné, nemusí se hlásit ÚOOÚ, ale porušení se má interně zdokumentovat.

Při hodnocení rizika se řeší povaha údajů, rozsah, identifikovatelnost, zranitelnost osob, zabezpečení, možnost zneužití a následky. Únik hashovaných hesel bude jiný, pokud jsou hesla silně solená a hashovaná moderním algoritmem, než pokud jde o plaintext hesla se jmény a e-maily. Pokud je organizace zároveň regulovaná podle ZKB, může vedle GDPR běžet i samostatné hlášení kybernetického incidentu.

Subjekty údajů mají práva podle čl. 12-22 GDPR: informace, přístup, opravu, výmaz, omezení zpracování, přenositelnost, námitku a nebýt předmětem čistě automatizovaného rozhodování s právními nebo obdobně významnými účinky. Žádost se vyřizuje zpravidla do jednoho měsíce, s možností prodloužení o dva měsíce u složitých případů.

\subsubsection*{Námitka vs výmaz}
\textbf{Námitka} podle čl. 21 míří hlavně na zpracování založené na oprávněném zájmu nebo veřejném zájmu. U přímého marketingu je zvlášť silná: pokud subjekt vznese námitku proti přímému marketingu, údaje už pro tento účel nesmí být zpracovávány. U jiného oprávněného zájmu může správce pokračovat jen tehdy, pokud prokáže závažné oprávněné důvody, které převažují nad zájmy a právy subjektu, nebo pokud data potřebuje pro právní nároky.

\textbf{Výmaz} podle čl. 17 znamená odstranění údajů, ale není absolutní. Neuplatní se, pokud správce údaje potřebuje například ke splnění právní povinnosti, pro výkon nebo obhajobu právních nároků, veřejný zájem, archivaci nebo svobodu projevu. Prakticky: zákazník může požádat o výmaz marketingového profilu, ale faktury se nevymažou, pokud je správce musí držet podle účetních a daňových předpisů.

\subsection*{ÚOOÚ}
ÚOOÚ je nezávislý dozorový úřad. Kontroluje dodržování GDPR a českého zákona, řeší stížnosti, ukládá opatření a pokuty, schvaluje některé nástroje, konzultuje DPIA, vydává metodiku a spolupracuje v Evropském sboru pro ochranu osobních údajů (EDPB). Nezávislost je klíčová, protože úřad dohlíží i na veřejnou moc.

Sankce mohou dosáhnout 20 mil. eur nebo 4 procent celosvětového ročního obratu u závažnějších porušení, případně 10 mil. eur nebo 2 procenta u méně závažných porušení. V přeshraničních věcech funguje mechanismus one-stop-shop: hlavní dozorový úřad je typicky úřad státu hlavní provozovny správce, ostatní dotčené úřady spolupracují a EDPB pomáhá sjednocovat výklad.

\subsection*{Předávání mimo EHP}
Předání osobních údajů mimo Evropský hospodářský prostor je možné jen při zajištění srovnatelné ochrany. Nejjednodušší je rozhodnutí Komise o odpovídající ochraně. Pokud není, používají se standardní smluvní doložky, závazná podniková pravidla nebo výjimky pro specifické situace. Po Schrems II je potřeba u standardních doložek posoudit i právo cílové země a praktická rizika přístupu veřejné moci, tedy transfer impact assessment.

Pro USA existuje EU-US Data Privacy Framework jako rozhodnutí o odpovídající ochraně pro certifikované subjekty. K 16. 6. 2026 platí, ale po rozhodnutí Tribunálu ve věci Latombe z 3. 9. 2025 běží odvolací rovina u Soudního dvora, takže u zkoušky je dobré říct: rámec se používá, ale politicko-právní nejistota kolem transatlantických přenosů trvá.

\subsection*{Modelové situace}
\begin{enumerate}
  \item \textbf{E-shop}: doručení objednávky je plnění smlouvy, faktury právní povinnost, bezpečnostní logy oprávněný zájem, newsletter novému zájemci souhlas, nabídka obdobného zboží stávajícímu zákazníkovi může spadat pod soft opt-in. Námitka proti marketingu zastaví marketing, ale nevymaže faktury.
  \item \textbf{Kamerový systém v obci}: veřejná moc nebo oprávněný zájem podle konkrétního režimu, jasný účel, informační cedule, omezená doba uložení, záběr jen potřebných míst, případně DPIA při rozsáhlém systematickém monitorování.
  \item \textbf{HR screening uchazečů}: nelze bez dalšího sbírat vše z internetu. Účel musí souviset s pozicí, rozsah musí být přiměřený, uchazeč má být informován a zvláštní kategorie údajů jsou vysoce rizikové.
  \item \textbf{AI profilování klientů banky}: profilování kvůli prevenci podvodů může mít oprávněný zájem nebo právní povinnost, ale je potřeba posoudit rizika, vysvětlitelnost, lidský dohled, minimalizaci a často DPIA.
\end{enumerate}

\priklad{Avast/Jumpshot: ÚOOÚ uložil Avast Software pokutu 351 mil. Kč za neoprávněné zpracování údajů o prohlížení uživatelů. Klíčová učební pointa je pseudo vs anonym: odstranění přímých identifikátorů a nahrazení identifikátorem nemusí být anonymizace, pokud lze osoby rozumně znovu identifikovat nebo profilovat. Městský soud v Praze v roce 2025 potvrdil porušení, ale zrušil část rozhodnutí ohledně výše pokuty kvůli odůvodnění; věc tím dobře ukazuje, že "anonymizováno" nestačí jen tvrdit.}

\zkouska{U GDPR odpovědi nikdy nezačínej "potřebujeme souhlas". Začni účelem, pak právním titulem, zásadami, minimalizací, rizikem a dokumentací. U oprávněného zájmu řekni tři kroky LIA; u rizikového zpracování řekni DPIA před spuštěním; u citlivých údajů řekni čl. 6 plus čl. 9.}

\zdroje{GDPR 2016/679; zákon o zpracování osobních údajů (č. 110/2019 Sb.); zákon o některých službách informační společnosti (č. 480/2004 Sb.); metodiky ÚOOÚ k DPIA a DPO; EDPB Guidelines 1/2024 k oprávněnému zájmu; věc Breyer C-582/14; Schrems II C-311/18; věc Latombe T-553/23 a navazující odvolání; ÚOOÚ a Městský soud v Praze ve věci Avast/Jumpshot; přiložené PDF a vložený text.}

\newpage

\otazka{23. Elektronický podpis a elektronická pečeť}{Právní úprava a druhy. Datové schránky - právní úprava a praxe používání.}{BI301K}

\subsection*{Kostra odpovědi}
\begin{enumerate}
  \item Vysvětlit technický princip podpisu: hash, soukromý a veřejný klíč, certifikát.
  \item Vyložit eIDAS, eIDAS 2.0 a české zákony o službách vytvářejících důvěru, elektronické identifikaci a datových schránkách/autorizované konverzi.
  \item Popsat druhy podpisu: prostý, zaručený, kvalifikovaný; český uznávaný podpis.
  \item Popsat elektronickou pečeť a časové razítko.
  \item Vyložit datové schránky: zřízení, dodání, doručení, fikce, vztah k podpisu.
  \item Zmínit autorizovanou konverzi.
\end{enumerate}

\subsection*{Technický princip}
Elektronický podpis využívá asymetrickou kryptografii. Z dokumentu se vypočte hash, tedy krátký otisk. Podepisující použije soukromý klíč k vytvoření podpisové hodnoty. Ověřovatel použije veřejný klíč, ověří podpis a znovu vypočte hash dokumentu. Pokud se dokument změní, hash nesedí a podpis neprojde.

Veřejný klíč je spojen s osobou certifikátem. Certifikát vydává certifikační autorita. Ověření podpisu proto znamená nejen ověřit kryptografii, ale i řetězec certifikátů, platnost, účel certifikátu, revokaci a čas podpisu. Pro dlouhodobé ověřování je důležité časové razítko a uchování validačních dat.

\subsection*{Právní rámec}
Základem je eIDAS - nařízení EU 910/2014, novelizované eIDAS 2.0 nařízením 2024/1183. eIDAS upravuje elektronickou identifikaci, elektronické podpisy, pečetě, časová razítka, služby elektronického doporučeného doručování, certifikáty a důvěryhodné služby. Českou úpravu doplňuje zákon o službách vytvářejících důvěru pro elektronické transakce (č. 297/2016 Sb.), zákon o elektronické identifikaci (č. 250/2017 Sb.) a zákon o elektronických úkonech a autorizované konverzi dokumentů (č. 300/2008 Sb.).

eIDAS stanoví důležitou zásadu: elektronickému podpisu nesmí být upírány právní účinky jen proto, že je elektronický. Kvalifikovaný elektronický podpis má stejné právní účinky jako vlastnoruční podpis. V ČR od roku 2023 vykonává úkoly orgánu dohledu v oblasti služeb vytvářejících důvěru Digitální a informační agentura (DIA); dříve tuto agendu vykonávalo Ministerstvo vnitra. NÚKIB tedy není obecným dozorovým orgánem pro elektronické podpisy a pečetě.

eIDAS 2.0 přidává důraz na evropskou peněženku digitální identity (EUDI Wallet), atributy identity a praktičtější přeshraniční používání důvěryhodných služeb. Pro otázku na podpisy je pointa hlavně v tom, že se posiluje ekosystém elektronické identity a kvalifikovaných služeb: osoba se má umět důvěryhodně identifikovat, podepsat, prokazovat atributy a používat digitální doklady napříč EU. Peněženka nemá být jen "občanka v mobilu"; má umožnit online i offline prokazování identity, sdílení jen vybraných údajů, ukládání digitálních dokladů a vytváření právně závazných elektronických podpisů. Přesné načasování peněženek není dobré říkat jen jako "do konce 2026" bez kontextu: eIDAS 2.0 váže dostupnost peněženek na lhůty od vstupu relevantních prováděcích aktů v účinnost, takže před zkouškou je vhodné ověřit aktuální stav implementačních aktů a českého projektu DIA.

\pravobox{eIDAS 2.0 v praxi: každý členský stát má zpřístupnit alespoň jednu evropskou peněženku digitální identity občanům, rezidentům a podnikům. Uživatel bude moci prokázat identitu nebo atribut, například věk či oprávnění jednat za firmu, aniž by pokaždé sdílel celý doklad. V ČR přípravu EUDI Wallet koordinuje DIA. Navazuje to na český ekosystém elektronické identity, například Identitu občana, ale eDoklady nelze jednoduše označit za totéž co EUDI peněženku - jde o samostatné národní řešení a jeho vztah k evropské peněžence se postupně formuje.}

Pro podpisy je prakticky důležité, že peněženka má umožnit i vytváření elektronických podpisů. U fyzických osob se počítá s tím, že kvalifikovaný elektronický podpis pro neprofesní použití bude dostupný bez poplatku; neznamená to ale, že všechny podpisové služby pro firmy nebo profesní použití budou zdarma. eIDAS 2.0 také počítá s tím, že vybrané veřejné a soukromé služby budou muset peněženku přijímat, aby ji člověk mohl použít nejen vůči státu, ale i v běžných přeshraničních situacích, například u banky, telekomunikační služby nebo při prokazování kvalifikace.

\subsection*{Druhy elektronického podpisu}
\begin{itemize}
  \item \textbf{Prostý elektronický podpis} - jakákoli elektronická data připojená k projevu vůle: jméno v e-mailu, zaškrtnutí, sken podpisu. Má právní účinky, ale slabší důkazní jistotu.
  \item \textbf{Zaručený elektronický podpis} - splňuje požadavky vazby na podepisujícího, identifikace, výhradní kontroly a detekce změn.
  \item \textbf{Kvalifikovaný elektronický podpis (QES)} - zaručený podpis založený na kvalifikovaném certifikátu a vytvořený kvalifikovaným prostředkem. Má účinky vlastnoručního podpisu v celé EU.
\end{itemize}

Český pojem \textbf{uznávaný elektronický podpis} zahrnuje kvalifikovaný podpis a zaručený podpis založený na kvalifikovaném certifikátu. Není totožný s QES. Každý QES je uznávaný, ale ne každý uznávaný podpis je QES.

\pozor{Nepleť zaručený, kvalifikovaný a uznávaný podpis. "Uznávaný" je český pojem, "kvalifikovaný" je eIDAS pojem s účinky vlastnoručního podpisu.}

\pravobox{Rychlá pomůcka: \textbf{prostý podpis} může být skoro cokoli elektronického, ale důkazně je slabý. \textbf{Zaručený podpis} už musí technicky identifikovat podepisujícího a chránit integritu. \textbf{Kvalifikovaný podpis} je zaručený podpis na kvalifikovaném certifikátu a s kvalifikovaným prostředkem, proto má účinky vlastnoručního podpisu. \textbf{Uznávaný podpis} je česká procesní zkratka pro komunikaci s veřejnou mocí; zahrnuje QES, ale není s ním totožný.}

\subsection*{Elektronická pečeť a časové razítko}
Elektronická pečeť je obdobou podpisu pro právnickou osobu nebo instituci, ale není to totéž co podpis. \textbf{Elektronický podpis} náleží fyzické osobě a typicky zachycuje její projev vůle. \textbf{Elektronická pečeť} náleží právnické osobě a dokládá původ a integritu dokumentu od této organizace. I pečeť může být prostá, zaručená nebo kvalifikovaná. Kvalifikovaná pečeť zakládá silnou domněnku původu a integrity, ale není ekvivalentem vlastnoručního podpisu fyzické osoby a sama o sobě neprokazuje, že konkrétní člověk učinil právní jednání.

Elektronické časové razítko spojuje data s konkrétním časem a prokazuje, že data v dané podobě v daném čase existovala. Kvalifikované časové razítko má vyšší důkazní sílu. Dlouhodobá validace kombinuje podpis, certifikát, revokační informace a časová razítka tak, aby bylo možné podpis ověřit i po vypršení certifikátu nebo zestárnutí algoritmů.

Podpis, pečeť ani časové razítko nedokazují pravdivost obsahu. Dokazují původ, integritu a čas. Dokument může být formálně správně podepsaný a přesto obsahovat nepravdu.

U dlouhodobého uchování je důležité neponechat dokument jen jako "soubor s podpisem". Je potřeba zachovat validační data: certifikát, řetězec důvěry, informace o revokaci a časová razítka. V praxi se používají formáty a profily jako PAdES u PDF, XAdES u XML nebo ASiC kontejnery. Dlouhodobá validace ale neopraví podpis, který byl neplatný už v okamžiku vytvoření; jen umožňuje později doložit, že tehdy byl ověřitelný.

\subsection*{Služby vytvářející důvěru}
Podpisy, pečetě a razítka patří mezi služby vytvářející důvěru. Důležité je, že kvalifikovaný status se vztahuje vždy ke konkrétní službě, ne automaticky k celému poskytovateli. Poskytovatel může mít kvalifikovanou službu vydávání certifikátů pro podpisy, ale nemusí mít kvalifikovaná časová razítka nebo kvalifikovanou službu elektronického doporučeného doručování. Proto se v praxi ověřuje aktuální důvěryhodný seznam, místo aby se jen pamatoval název firmy.

Kvalifikované elektronické doporučené doručování podle eIDAS je evropská kategorie důvěryhodné služby. Umí prokazovat odeslání, přijetí, čas, integritu a identifikaci stran. Není to totéž co česká datová schránka. Funkčně se podobají tím, že obě řeší prokazatelné doručování, ale mají jiný právní základ a jinou územní logiku: datové schránky jsou český národní systém, kvalifikované doporučené doručování je evropská služba podle eIDAS.

\subsection*{Datové schránky}
Datové schránky upravuje zákon o elektronických úkonech a autorizované konverzi dokumentů. Slouží k elektronickému doručování a k činění úkonů vůči orgánům veřejné moci. Informační systém datových schránek spravuje DIA a provozuje jej Česká pošta. Datová zpráva je dodána do schránky okamžikem dodání a doručena přihlášením oprávněné osoby. Fikce doručení po 10 dnech je typická hlavně u doručování orgánem veřejné moci; u soukromé komunikace a poštovních datových zpráv je potřeba rozlišit konkrétní režim a neříkat ji úplně univerzálně.

Informační systém datových schránek je veřejnoprávní doručovací infrastruktura. Pro zkoušku je důležité rozlišovat \textbf{dodání} a \textbf{doručení}. Dodání znamená, že zpráva technicky dorazila do schránky. Doručení znamená právní okamžik, kdy se s ní adresát má seznámit - typicky přihlášením oprávněné osoby nebo fikcí po 10 dnech. U orgánů veřejné moci a soukromých osob se liší praktické režimy doručování a dostupnost některých služeb, ale základní pointa je stejná: datová schránka je důvěryhodný kanál, ne podpisový nástroj.

Je dobré rozlišit tři komunikační situace. \textbf{Orgán veřejné moci} doručuje do datové schránky v režimu veřejnoprávního doručování. \textbf{Soukromá osoba vůči orgánu veřejné moci} typicky činí podání, kterému zákon přiznává účinky podepsaného úkonu. \textbf{Soukromá komunikace mezi soukromými osobami} funguje jako poštovní datová zpráva nebo soukromoprávní komunikace, kde se právní následky a fikce neposuzují úplně stejně jako u doručování orgánu veřejné moci. Proto je potřeba vždy říct, kdo komu doručuje a v jakém režimu.

Datová schránka není elektronický podpis. Odeslání z datové schránky ale může mít účinky podepsaného úkonu, pokud to zákon připouští. Jde o zákonné nahrazení podpisu, ne o vytvoření podpisu uvnitř dokumentu. Proto je nutné rozlišit:
\begin{itemize}
  \item dokument podepsaný elektronickým podpisem,
  \item dokument odeslaný datovou schránkou,
  \item dokument autorizovaně konvertovaný,
  \item běžný sken vlastnoručně podepsaného dokumentu.
\end{itemize}

Datové zprávy nejsou ve schránce uchovávány navždy. V praxi je nutné řešit archivaci, doručenky a dlouhodobé ověření dokumentů. Pro právní jistotu je vhodné ukládat nejen přílohu, ale i obálku datové zprávy a doručenku.

Zřízení datové schránky se liší podle typu osoby. Orgány veřejné moci, právnické osoby zapsané v rejstřících, podnikající fyzické osoby a vybrané profese ji mají typicky zřízenou ze zákona; nepodnikající fyzická osoba ji má na žádost. Prakticky je důležité vědět i to, že ISDS není archiv. Zprávy se uchovávají jen omezenou dobu, řádově 90 dnů po doručení, a pro dlouhodobé uchování je potřeba spisová služba, Datový trezor, stažení celé zprávy včetně doručenky nebo autorizovaná konverze.

\pozor{Datová schránka často prakticky nahradí podpis podání, ale nevytvoří elektronický podpis uvnitř PDF. Když později oddělíš přílohu od obálky datové zprávy a doručenky, můžeš přijít o důkazní kontext, kdy a odkud byla zpráva odeslána a doručena.}

\subsection*{Autorizovaná konverze}
Autorizovaná konverze převádí dokument mezi listinnou a elektronickou podobou se zachováním právních účinků formy. Není to obyčejný sken, tisk ani ověření pravdivosti obsahu. Konverze potvrzuje shodu podoby dokumentu. Při konverzi elektronického dokumentu se kontrolují elektronické prvky, například podpisy, pečetě a časová razítka.

Výstup konverze obsahuje konverzní doložku. Ta říká, kdo, kdy a z čeho dokument převedl, a potvrzuje shodu podoby vstupu a výstupu. Neznamená to, že obsah dokumentu je pravdivý nebo že právní jednání v něm je platné. Konverze řeší formu a důkazní použitelnost, ne věcnou správnost.

\priklad{Podání soudu lze učinit datovou schránkou; podpis se pak zpravidla nevyžaduje, protože zákon přiznává úkonu účinky podepsaného podání. Když stejný dokument pošlu obyčejným e-mailem bez uznávaného podpisu, účinky mohou být jiné.}

\zkouska{Závěrečná věta může být: "Elektronický podpis řeší, kdo a s jakou integritou dokument podepsal; datová schránka řeší důvěryhodný kanál doručení a identifikaci schránky; autorizovaná konverze řeší převod formy dokumentu."}

\zdroje{eIDAS 910/2014 a eIDAS 2.0 - nařízení 2024/1183; zákon o službách vytvářejících důvěru pro elektronické transakce (č. 297/2016 Sb.); zákon o elektronické identifikaci (č. 250/2017 Sb.); zákon o elektronických úkonech a autorizované konverzi dokumentů (č. 300/2008 Sb.); přiložené PDF.}

\newpage
\section*{Rychlé opakování před zkouškou}
\addcontentsline{toc}{section}{Rychlé opakování před zkouškou}
\begin{itemize}
  \item Hrozba není riziko. Hrozba je zdroj škody, riziko je míra ohrožení konkrétního objektu.
  \item Kyberbezpečnost stírá hranici vnitřní a vnější bezpečnosti.
  \item Strategie není zákon. Strategie určuje cíle, zákon ukládá povinnosti, akční plán rozepisuje úkoly.
  \item Ne každý kyberútok je kyberválka. Většina operací je pod prahem ozbrojeného konfliktu.
  \item U kritické infrastruktury nejde jen o data, ale o kontinuitu základní služby a fyzické dopady.
  \item NIS2 se v ČR učí přes zákon o kybernetické bezpečnosti a prováděcí vyhlášky.
  \item U kyberkriminality převáděj technický popis na skutkové podstaty.
  \item U elektronických důkazů hlídej volatilitu, integritu, hash a chain of custody.
  \item U GDPR nejdřív účel, potom právní titul, minimalizace, riziko a dokumentace.
  \item Datová schránka není elektronický podpis, ale zákon jí může dát účinky podepsaného úkonu.
\end{itemize}

\section*{Dynamické údaje k ověření před státnicí}
\addcontentsline{toc}{section}{Dynamické údaje k ověření před státnicí}
\begin{itemize}
  \item \textbf{E-evidence}: od 18. 8. 2026 se začne používat EPOC/EPOC-PR; před zkouškou ověř, zda už je režim použitelný, jak je nastaven český příjem a validace příkazů a zda se změnila praxe poskytovatelů.
  \item \textbf{2. dodatkový protokol k Budapešťské úmluvě}: ověř počet ratifikací a zda už vstoupil v platnost; k 16. 6. 2026 mu chyběla pátá ratifikace.
  \item \textbf{Úmluva OSN proti kyberkriminalitě}: ověř počet signatářů, ratifikací a vstup v platnost; k 16. 6. 2026 nebyla v platnosti.
  \item \textbf{CER/ZKI}: ověř prováděcí nařízení k základním službám a významovým kritériím, přechodné lhůty a to, zda už byly určeny nové kritické subjekty.
  \item \textbf{ZKB/NIS2}: ověř aktuální prováděcí předpisy NÚKIB, metodiky, režim strategicky významných služeb a pravidla bezpečnosti dodavatelského řetězce.
  \item \textbf{eIDAS 2.0 a EUDI Wallet}: ověř prováděcí akty, český projekt DIA, povinnosti přijímat peněženku a přesné lhůty pro zpřístupnění.
  \item \textbf{Předávání osobních údajů do USA}: ověř stav EU-US Data Privacy Framework a případná rozhodnutí Soudního dvora EU navazující na věc Latombe.
  \item \textbf{Strategie}: ověř aktuální akční plán Národní strategie kybernetické bezpečnosti 2026-2030 a dokumenty ke Strategii pro posílení odolnosti subjektů kritické infrastruktury 2026-2029.
\end{itemize}

\section*{Ověřené online zdroje}
\addcontentsline{toc}{section}{Ověřené online zdroje}
\begin{itemize}
  \item NÚKIB - strategie a akční plány: \url{https://nukib.gov.cz/cs/infoservis/dokumenty-a-publikace/strategie-akcni-plan/}
  \item Portál NÚKIB k novému zákonu o kybernetické bezpečnosti: \url{https://portal.nukib.gov.cz/}
  \item e-Sbírka - zákon o kybernetické bezpečnosti (č. 264/2025 Sb.): \url{https://e-sbirka.gov.cz/sb/2025/264}
  \item e-Sbírka - zákon o odolnosti subjektů kritické infrastruktury (č. 266/2025 Sb.): \url{https://e-sbirka.gov.cz/sb/2025/266}
  \item HZS ČR - kritická infrastruktura po nové úpravě: \url{https://hzscr.gov.cz/clanek/web-informacni-servis-zpravodajstvi-2026-unor-kriticka-infrastruktura.aspx}
  \item HZS ČR - představení zákona o kritické infrastruktuře: \url{https://hzscr.gov.cz/clanek/predstaveni-zakona-o-kriticke-infrastrukture.aspx}
  \item Ministerstvo vnitra - poskytování informací podle zákona o kritické infrastruktuře: \url{https://mv.gov.cz/npo/clanek/sdeleni-ministerstva-vnitra-k-poskytovani-informaci-podle-zakona-o-kriticke-infrastrukture.aspx}
  \item Vláda ČR - Bezpečnostní rada státu: \url{https://vlada.gov.cz/cz/pracovni-a-poradni-organy-vlady/brs/brs-uvod-3851/}
  \item EUR-Lex - NIS2: \url{https://eur-lex.europa.eu/eli/dir/2022/2555/oj}
  \item EUR-Lex - CER: \url{https://eur-lex.europa.eu/eli/dir/2022/2557/oj}
  \item EUR-Lex - Cybersecurity Act 2019/881: \url{https://eur-lex.europa.eu/eli/reg/2019/881/oj}
  \item EUR-Lex - Cyber Solidarity Act 2025/38: \url{https://eur-lex.europa.eu/eli/reg/2025/38/oj}
  \item EUR-Lex - útoky na informační systémy, směrnice 2013/40/EU: \url{https://eur-lex.europa.eu/eli/dir/2013/40/oj}
  \item EUR-Lex - bezhotovostní platební podvody, směrnice 2019/713: \url{https://eur-lex.europa.eu/eli/dir/2019/713/oj}
  \item EUR-Lex - evropský vyšetřovací příkaz, směrnice 2014/41/EU: \url{https://eur-lex.europa.eu/eli/dir/2014/41/oj}
  \item EUR-Lex - e-evidence nařízení 2023/1543: \url{https://eur-lex.europa.eu/eli/reg/2023/1543/oj}
  \item EUR-Lex - e-evidence směrnice 2023/1544: \url{https://eur-lex.europa.eu/eli/dir/2023/1544/oj}
  \item EUR-Lex - Cyber Resilience Act 2024/2847: \url{https://eur-lex.europa.eu/eli/reg/2024/2847/oj}
  \item EUR-Lex - eIDAS 2.0, nařízení 2024/1183: \url{https://eur-lex.europa.eu/eli/reg/2024/1183/oj}
  \item EUR-Lex - GDPR, nařízení 2016/679: \url{https://eur-lex.europa.eu/eli/reg/2016/679/oj}
  \item EUR-Lex - policejní směrnice 2016/680: \url{https://eur-lex.europa.eu/eli/dir/2016/680/oj}
  \item ÚOOÚ - posouzení vlivu na ochranu osobních údajů: \url{https://uoou.gov.cz/profesional/metodiky-a-doporuceni-pro-spravce/posouzeni-vlivu-na-ochranu-osobnich-udaju}
  \item ÚOOÚ - základní příručka k ochraně údajů: \url{https://uoou.gov.cz/verejnost/zakladni-prirucka-k-ochrane-udaju}
  \item EDPB - územní působnost GDPR: \url{https://www.edpb.europa.eu/our-work-tools/our-documents/guidelines/guidelines-32018-territorial-scope-gdpr-article-3-version_en}
  \item EDPB - Guidelines 1/2024 k oprávněnému zájmu: \url{https://www.edpb.europa.eu/our-work-tools/documents/public-consultations/2024/guidelines-12024-processing-personal-data-based_en}
  \item ÚOOÚ - Avast/Jumpshot pokuta 351 mil. Kč: \url{https://uoou.gov.cz/novinky/vse/uoou-ulozil-pokutu-351-mil-kc-za-poruseni-gdpr}
  \item Městský soud v Praze - Avast/Jumpshot rozsudek 2025: \url{https://msp.gov.cz/documents/d/mestsky-soud-v-praze/tiskova-zprava-k-rozsudku-ze-dne-30-9-2025-sp-zn-5-a-56-2024}
  \item CURIA - Latombe v Commission, T-553/23: \url{https://infocuria.curia.europa.eu/tabs/redirect/juris/liste.jsf?num=T-553/23}
  \item Council of Europe - Budapešťská úmluva a protokoly: \url{https://www.coe.int/en/web/cybercrime/the-budapest-convention}
  \item Council of Europe - 2. dodatkový protokol CETS 224, podpisy a ratifikace: \url{https://www.coe.int/en/web/conventions/full-list?module=signatures-by-treaty&treatynum=224}
  \item UN Treaty Collection - Úmluva OSN proti kyberkriminalitě: \url{https://treaties.un.org/pages/ViewDetails.aspx?chapter=18&clang=_en&mtdsg_no=XVIII-16&src=TREATY}
  \item DIA - evropská peněženka digitální identity: \url{https://www.dia.gov.cz/eudiw/cs}
  \item NATO - kolektivní obrana a článek 5: \url{https://www.nato.int/en/what-we-do/introduction-to-nato/collective-defence-and-article-5}
\end{itemize}

\end{document}
